\BourbakiExerciseGroup

\begin{enumerate}
\def\labelenumi{\arabic{enumi})}
\item
  Let E be an extension of a field K and x and y two distinct roots in E of the same irreducible polynomial of K{[}X{]}. Show that the extensions \(\mathbf{K}(x)\) and \(\mathbf{K}(y)\) are not linearly disjoint over K (use Th. 2). * Give an example where K = Q and \(\mathbf{K}(x) \cap \mathbf{K}(y) = \mathbf{K}\) (cf.~V, p.~161, Ex. 2). *
\item
  Let \((\mathbf{E}_i)_{i\in I}\) be a family of extensions of a field \(\mathbf{K}\) , contained in an extension \(\mathbf{L}\) of \(\mathbf{K}\) . If \(\mathbf{F}_i\) is the relative algebraic closure of \(\mathbf{K}\) in \(\mathbf{E}_i\) , show that the relative algebraic closure of \(\mathbf{K}\) in \(\mathbf{E} = \bigcap_{i\in I}\mathbf{E}_i\) is \(\mathbf{F} = \bigcap_{i\in I}\mathbf{F}_i\) .
\item
  Show that every algebraic extension E of a field K is equipotent to a subset of \(K \times N\) (consider the mapping which associates with each element of E its minimal polynomial over K). In particular, every algebraic extension of a finite field is countable and every algebraic extension of an infinite field K is equipotent to K. * Deduce that in the field R of real numbers there exist numbers that are transcendental over the prime subfield Q and that the set of all these numbers has the power of the continuum. *
\item
  Let F be a transcendental extension of a field E. Show that
\end{enumerate}

\[
[ \mathrm{F}: \mathrm{E} ] \geqslant \operatorname * {S u p} (\operatorname{Card} (\mathrm{E}), \operatorname{Card} (\mathbf {N}))
\]

with equality if the extension F can be generated by a countable family of elements. (Use Ex. 7 of IV, p.~89.)

\begin{enumerate}
\def\labelenumi{\arabic{enumi})}
\setcounter{enumi}{4}
\tightlist
\item
  Let A be a commutative algebra and let S be a subset of A generating A as K-algebra. Suppose that \(\operatorname{Card}(S) < \operatorname{Card}(K)\) . Show that if M is a maximal ideal of A, then the field A/M is an algebraic extension of K. (Use Ex. 15 of V, p.174 when \(\operatorname{Card}(K) \leqslant \operatorname{Card}(N)\) and Ex. 4 when
\end{enumerate}

\[
\operatorname{Card} (\mathbf {K}) > \operatorname{Card} (\mathbf {N}).
\]

* Deduce that \(\mathbf{A} / \mathfrak{M} = \mathbf{K}\) when \(\mathbf{K}\) is algebraically closed. *

* 6) Let C be an uncountable algebraically closed field and \(\mathbf{P} = \mathbf{C}[(\mathbf{X}_i)_{i \in I}]\) a polynomial algebra over C. Given two families \((f_\alpha)_{\alpha \in A}, (g_\beta)_{\beta \in B}\) of elements of P suppose that: a) B and I are countable.

\begin{enumerate}
\def\labelenumi{\alph{enumi})}
\setcounter{enumi}{1}
\tightlist
\item
  For every finite subset \(\mathbf{A}'\) of \(\mathbf{A}\) and every finite subset \(\mathbf{B}'\) of \(\mathbf{B}\) there exists \((x_i)_{i\in I} \in C^I\) such that
\end{enumerate}

\[
f _ {\alpha} (x _ {i}) = 0 \quad \text { and } \quad g _ {\beta} (x _ {i}) \neq 0 \quad \text { for   all } \quad \alpha \in \mathbf {A} ^ {\prime}, \beta \in \mathbf {B} ^ {\prime}.
\]

Show that then there exists \((x_{i})_{i\in I}\in \mathbf{C}^{I}\) such that

\[
f _ {\alpha} (x _ {i}) = 0 \quad \text { and } \quad g _ {\beta} (x _ {i}) \neq 0 \quad \text { for   all } \quad \alpha \in A, \beta \in B.
\]

(Let \((\mathbf{Y}_{\beta})_{\beta \in \mathbb{B}}\) be a family of indeterminates indexed by \(\mathbf{B}\) and let \(\mathbf{Q}\) be the polynomial algebra \(\mathbf{C}[(\mathbf{X}_i)_{i\in I}, (\mathbf{Y}_{\beta})_{\beta \in \mathbb{B}}]\) . Let \(\Lambda\) be the quotient of \(\mathbf{Q}\) by the ideal generated by the \(f_{\alpha}\) and the \(1 - g_{\alpha}\mathbf{Y}_{\alpha}\) . Show that \(\Lambda \neq 0\) thanks to \(a)\) and \(b)\) and apply Ex. 5 to a maximal ideal of \(\Lambda\) .) *

\begin{enumerate}
\def\labelenumi{\arabic{enumi})}
\setcounter{enumi}{6}
\tightlist
\item
  Let \(L\) be a field, \(A\) a subring of \(L\) and \(K \subset L\) the field of fractions of \(A\) .
\end{enumerate}

\begin{enumerate}
\def\labelenumi{\alph{enumi})}
\item
  Show that if L is a finitely generated A-module (II, p.~206) then necessarily A = K (by considering a supplementary subspace of K in L, qua vector space over K, show that K is a finitely generated A-module, hence necessarily of the form \(s^{-1}A\) , where \(s \in A\) ; finally show that s must be invertible in A).
\item
  Show that if there exist a finite number of elements \(x_{j}\) ( \(1 \leqslant j \leqslant n\) ) of L, algebraic over K, such that \(L = A[x_{1}, \ldots, x_{n}]\) , then there exists an element \(b \neq 0\) in A such that \(K = A[b^{-1}]\) (show that there exists \(b \neq 0\) in A such that L is a finitely generated \(A[b^{-1}]\) -module). Show that either b is invertible (and K = A) or every non-zero ideal of A contains a principal ideal \(Ab^{m}\) . Converse. The ring \(k[[X]]\) of formal power series over a field k is a ring of this type.
\end{enumerate}

* 8) If K is a field, show that K is algebraically closed in the field K(X) of rational fractions in one indeterminate over K (use the fact that K{[}X{]} is a principal ideal domain). *

\begin{enumerate}
\def\labelenumi{\arabic{enumi})}
\setcounter{enumi}{8}
\tightlist
\item
  Let \(\mathbf{K}\) be a field, \(\mathbf{E}\) an extension of \(\mathbf{K}\) and \(\mathbf{S}\) a subset of \(\mathbf{E} - \mathbf{K}\) .
\end{enumerate}

\begin{enumerate}
\def\labelenumi{\alph{enumi})}
\item
  Show that the set of extensions \(\mathbf{L} \subset \mathbf{E}\) of \(\mathbf{K}\) such that \(\mathbf{L} \cap \mathbf{S} = \varnothing\) , ordered by inclusion, has a maximal element \(\mathbf{M}\) .
\item
  Show that if S is finite, E is an algebraic extension of M.
\end{enumerate}

* 10) Let K be a field and P a polynomial in the ring K{[}X, Y{]} in two indeterminates. Suppose that there exists a polynomial H ∈ K{[}X{]} and two polynomials Q, R ∈ K{[}X, Y{]} such that H(X) P(X, Y) = Q(X, Y) R(X, Y) and that the coefficients of Q(X, Y), qua polynomial in Y, are polynomials in K{[}X{]} without non-constant common factor. Show that all the coefficients of R(X, Y), qua polynomial in Y, are divisible by H(X). (Consider H, P, Q, R as polynomials in X over the field K(Y) and decompose them into factors of the first degree in Ω{[}X{]} where Ω is an algebraic closure of K(Y); observe that if H(X) and Q(X, Y) had a factor of the first degree in common, this factor would divide all the coefficients of Q(X, Y), qua polynomial in Y, and use the fact (IV, p.~12) that if K' is an extension of K, the GCD in K'{[}X{]} of two polynomials of K{[}X{]} belongs to K{[}X{]}.) *

\begin{enumerate}
\def\labelenumi{\arabic{enumi})}
\setcounter{enumi}{10}
\tightlist
\item
  Let \(\mathbf{E}\) be an extension of a field \(\mathbf{K}, x \in \mathbf{E}\) a transcendental element over \(\mathbf{K}\) , so that \(\mathbf{K}(x)\) is isomorphic to the field \(\mathbf{K}(\mathbf{T})\) of rational fractions in an indeterminate over \(\mathbf{K}\) . Thus every \(y \in \mathbf{K}(x)\) may be written \(y = g(x) / h(x)\) , where \(g\) and \(h\) are two polynomials of \(\mathbf{K}[\mathbf{X}]\) that are relatively prime (IV, p.~12), determined up to a non-zero common factor in \(\mathbf{K}\) ; the larger of the degrees of \(g\) and \(h\) is called the height of \(y\) with respect to \(x\) .
\end{enumerate}

\begin{enumerate}
\def\labelenumi{\alph{enumi})}
\tightlist
\item
  Show that in the polynomial ring \(\mathbf{K}(y)[\mathbf{X}]\) the polynomial \(g(\mathbf{X}) - yh(\mathbf{X})\) is irreducible if \(y \notin \mathbf{K}\) . Deduce that if \(y\) has height \(n > 0\) with respect to \(x\) , then \(\mathbf{K}(x)\) is an algebraic extension of degree \(n\) of \(\mathbf{K}(y)\) .
\end{enumerate}

Let \(\mathrm{P}(\mathbf{X}) = \sum_{i=0}^{n} y_i \mathbf{X}^i\) be the minimal polynomial of \(x\) over \(\mathrm{K}(y)\) . Show that for

\(0 \leqslant i \leqslant n\) we have either \(y_i \in \mathbf{K}\) or \(\mathbf{K}(y_i) = \mathbf{K}(y)\) .

\begin{enumerate}
\def\labelenumi{\alph{enumi})}
\setcounter{enumi}{1}
\item
  Deduce from a) that every \(y \in \mathbf{K}(x)\) such that \(\mathbf{K}(y) = \mathbf{K}(x)\) is of the form \((ax + b)/(cx + d)\) where a, b, c, d are elements of K such that \(ad - bc \neq 0\) ; converse. Find all K-automorphisms of \(\mathbf{K}(x)\) .
\item
  Show that if \(y \in \mathbf{K}(x)\) is of height \(n\) with respect to \(x\) , and \(z \in \mathbf{K}(y)\) of height \(m\) with respect to \(y\) , then \(z\) is of height \(mn\) with respect to \(x\) .
\end{enumerate}

* \(d\) ) Let F be an extension of K such that \(\mathbf{K} \subset \mathbf{F} \subset \mathbf{K}(x)\) and \(\mathbf{F} \neq \mathbf{K}\) ; let y be an element of F whose height \(m\) with respect to \(x\) has the least possible value; show that \(\mathbf{F} = \mathbf{K}(y)\) («Lüroth's theorem »). (Let P be the minimal polynomial of \(x\) in F{[}T{]}; show that \(\mathrm{P(T)} = \mathrm{Q(T,x)}/\mathrm{R}(x)\) , where Q is a polynomial of K{[}T,X{]} of degree \(\geqslant m\) in X and is not divisible by any non-constant polynomial of K{[}X{]}, and R is a polynomial of K{[}X{]}. If \(y = g(x)/h(x)\) , where \(g\) and \(h\) are relatively prime in K{[}X{]}, observe that by \(a\) ) the polynomial \(g(T)h(X) - g(X)h(T)\) is not divisible by any non-constant polynomial of K{[}T{]} or of K{[}X{]}; deduce that necessarily

\[
g (T) h (X) - g (X) h (T) = c. Q (T, X), \quad \text { where } c \in K,
\]

using Ex. 10.) *

* 12) a) With the notation and hypotheses as in Ex. 11, show that for an extension F of K such that \(\mathbf{K} \subset \mathbf{F} \subset \mathbf{K}(x)\) to be such that \(\mathbf{F} \cap \mathbf{K}[x] \neq \mathbf{K}\) , it is necessary and sufficient that \(\mathbf{F} = \mathbf{K}(y)\) where \(y = g(x)\) for a non-constant polynomial \(g \in \mathbf{K}[\mathrm{T}]\) ; then we have \(\mathbf{K}(y) \cap \mathbf{K}[x] = \mathbf{K}[y]\) . (Write \(\mathbf{F} = \mathbf{K}(y)\) with \(y = g(x)/h(x)\) , \(g\) and \(h\) being relatively prime polynomials of \(\mathbf{K}[\mathrm{T}]\) . If there exists a non-constant polynomial \(\mathbf{P} \in \mathbf{K}[\mathrm{T}]\) such that \(\mathbf{P}(x) \in \mathbf{F}\) , we can write \(\mathbf{P}(x) = \mathbf{Q}(y)/\mathbf{R}(y)\) , where \(\mathbf{Q}\) and \(\mathbf{R}\) are two relatively prime monic polynomials of \(\mathbf{K}[\mathrm{T}]\) , not both constant. Distinguish two cases, according as \(\mathbf{R} = 1\) or \(\mathbf{R}\) is non-constant; in the second case, decompose \(\mathbf{Q}\) and \(\mathbf{R}\) into factors of the first degree in an algebraic closure \(\Omega\) of \(\mathbf{K}\) and observe that if \(a, b\) are distinct elements of \(\Omega\) , then \(g - ah\) and \(g - bh\) are relatively prime in \(\Omega[\mathrm{T}]\) .)

\begin{enumerate}
\def\labelenumi{\alph{enumi})}
\setcounter{enumi}{1}
\item
  If \(\mathbf{F} = \mathbf{K}(y)\) with \(y = g(x)\) , where \(g \in \mathbf{K}[T]\) is non-constant, we may suppose that \(g(T) = Tg_1(T)\) , where \(g_1 \in \mathbf{K}[T]\) . Show that if \(\mathbf{K}(xg(x)) \cap \mathbf{K}[y] \neq \mathbf{K}\) , we necessarily have \(g_1(T) = aT^n\) with \(a \in \mathbf{K}\) and \(n > 0\) . (Observe that by hypothesis and \(a\) ) there exist two non-constant polynomials \(P, Q\) in \(K[T]\) such that \(P(xg(x)) = Q(g(x)) \notin K\) ; deduce that there are two elements \(a, b\) in \(K\) and an integer \(m\) such that \(g(T)\) divides \(a - bT^m\) .) c) Let \(y, z\) be two elements of \(K[x]\) such that \(K(y) \cap K(z) \neq K\) ; show that then \(K[y] \cap K[z] \neq K\) .
\item
  Find \(y \in \mathbf{K}(x)\) such that \(\mathbf{K}(y) = \mathbf{K}(x)\) but \(\mathbf{K}[y] \cap \mathbf{K}[x] = \mathbf{K}\) . *
\end{enumerate}

* 13) Let K be a field, \(P(T, X) = a_0(T)X^n + a_1(T)X^{n-1} + \cdots + a_n(T)\) a polynomial in K{[}T, X{]}; suppose that: \(1^\circ\) the polynomial \(a_0\) is not divisible by T; \(2^\circ\) each of the polynomials \(a_1, \ldots, a_n\) is divisible by T; \(3^\circ\) the polynomial \(a_n\) is not divisible by \(T^2\) . Show that under these conditions P, as polynomial of K(T){[}X{]}, is irreducible (argue by contradiction, using Ex. 10). *

* 14: With the notation as in Ex. 11, show that if we take \(g(X) = X\) and \(h(X) = X^n + X + 1 \quad (n \geqslant 2)\) , the polynomial \((g(X)h(x) - h(X)g(x)) / (X - x)\) of \(K(x)[X]\) is irreducible (use Ex. 13). *

* 15) Give an example of a field K containing two subfields \(\mathbf{K}_1\) , \(\mathbf{K}_2\) such that K is algebraic over \(\mathbf{K}_1\) and over \(\mathbf{K}_2\) but not over \(\mathbf{K}_1 \cap \mathbf{K}_2\) . (Take \(\mathbf{K} = \mathbf{Q}(T)\) , \(\mathbf{K}_1 = \mathbf{Q}(T^2)\) , and \(\mathbf{K}_2 = \mathbf{Q}(T^2 - T)\) . Verify that \([\mathbf{K} : \mathbf{K}_1] = 2\) , \([\mathbf{K} : \mathbf{K}_2] = 2\) and \(\mathbf{K}_1 \cap \mathbf{K}_2 = \mathbf{Q}\) . To do this note that an element \(f(T)\) of \(\mathbf{K}_1 \cap \mathbf{K}_2\) satisfies the conditions \(f(-T) = f(T)\) and \(f(1 - T) = f(T)\) .)

\BourbakiExerciseGroup

\begin{enumerate}
\def\labelenumi{\arabic{enumi})}
\tightlist
\item
  Let K be a field and \((\mathrm{E}_{i})_{i\in I}\) any family of extensions of K. Show that there exists an extension E of K and for each \(i\in I\) , a K-isomorphism \(u_{i}\) of \(E_{i}\) into E such that E is generated by the union of the \(u_{i}(\mathrm{E}_{i})\) (argue as in Prop. 4 (V, p.~22)).
\end{enumerate}

* 2) Let \(\mathbf{K}\) be an algebraically closed field of characteristic 0, \(\mathrm{F} = \mathrm{K}((\mathbf{X}^{1 / n!}))\) the field of formal power series in \(\mathbf{X}^{1 / n!}\) , and \(\mathrm{E} = \mathrm{K}((\mathbf{X}))\) , where \(\mathbf{X} = (\mathbf{X}^{1 / n!})^{n!}\) , the subfield of \(\mathbf{F}\) consisting of formal power series in \(\mathbf{X}\) . For every formal power series \(f \in \mathbf{K}((\mathbf{X}))\) we denote by \(\omega(f)\) its order.

\begin{enumerate}
\def\labelenumi{\alph{enumi})}
\item
  Let \(P(Y) = a_{0}Y^{n} + a_{1}Y^{n-1} + \cdots + a_{n}\) be a polynomial of E{[}Y{]} of degree n, such that \(a_{n} \neq 0\) . Show that there exists a strictly increasing sequence \((i_{k})_{0 \leqslant k \leqslant r}\) of integers in the interval \([0, n]\) such that: \(1^{\circ} i_{0} = 0\) , \(i_{r} = n\) ; \(2^{\circ} \omega(a_{i_{k}})\) is finite for \(0 \leqslant k \leqslant r\) ; \(3^{\circ}\) for each index j such that \(0 \leqslant j \leqslant n\) , distinct from the \(i_{k}\) and such that \(\omega(a_{j})\) is finite, the point \((j, \omega(a_{j})) \in \mathbf{R}^{2}\) is above the straight line passing through the points \((i_{k}, \omega(a_{i_{k}}))\) , \((i_{k-1}, \omega(a_{i_{k-1}}))\) and strictly above it if \(j < i_{k-1}\) or \(j > i_{k} (1 \leqslant k \leqslant r)\) . The union of the segments joining the points \((i_{k-1}, \omega(a_{i_{k-1}}))\) and \((i_{k}, \omega(a_{i_{k}}))\) for \(1 \leqslant k \leqslant r\) is called the Newton polygon of P, the preceding segments are called its sides and the points \((i_{k}, \omega(a_{i_{k}}))\) its vertices.
\item
  Put \(\rho_{k} = i_{k} - i_{k - 1}\) , \(\sigma_{k} = (\omega (a_{i_{k}}) - \omega (a_{i_{k - 1}})) / \rho_{k}\) ; show that \(P\) has exactly \(\rho_{k}\) zeros (counted with the right multiplicity) belonging to the field \(K((X^{1 / (\rho_k!)}))\) and whose orders are equal to \(\sigma_{k}\) (determine by recursion the coefficients of a formal power series of this type satisfying the equation \(P(z) = 0\) ).
\item
  Conclude that the union of the fields \(\mathbf{K}((\mathbf{X}^{1 / n}))\) for all integers \(n > 0\) is an algebraic closure of the field \(\mathrm{E} = \mathbf{K}((\mathbf{X}))\) (« Puiseux's theorem »).
\item
  If \(\mathbf{K}'\) is a field of characteristic 2 and \(\mathbf{E}' = \mathbf{K}'((\mathbf{X}))\) , then the polynomial \(\mathbf{Y}^2 + \mathbf{X}\mathbf{Y} + \mathbf{X}\) of \(\mathbf{E}'[\mathbf{Y}]\) does not have a root in the union of the fields \(\mathbf{K}'((\mathbf{X}^{1/n!}))\) . *
\end{enumerate}

\BourbakiExerciseGroup

\begin{enumerate}
\def\labelenumi{\arabic{enumi})}
\tightlist
\item
  \begin{enumerate}
  \def\labelenumii{\alph{enumii})}
  \tightlist
  \item
    Let E be an extension of a field K, \(x \in E\) a transcendental element over K and let \(F = K(x^{n})\) for an integer n \textgreater{} 1. Show that if the intersection \(F \cap K[(x - 1)(x^{n} - 1)]\) is distinct from K, then K is of characteristic p \textgreater{} 0 and \(n = p^{e}\) , so that \(K(x)\) is a p-radical extension of F. (Writing y = x - 1 and using Ex. 12, b) of V, p.~149, show that \((y + 1)^{n} - 1 = ay^{n}\) for an \(a \in K\) .)
  \end{enumerate}
\end{enumerate}

\begin{enumerate}
\def\labelenumi{\alph{enumi})}
\setcounter{enumi}{1}
\item
  Show that if \(\mathbf{K} \subset \mathbf{F}' \subset \mathbf{K}(x)\) and if \(\mathbf{K}(x)\) is not a \(p\) -radical extension of \(\mathbf{F}'\) and if \(\mathbf{K} \neq \mathbf{F}' \neq \mathbf{K}(x)\) , then there exists \(z \in \mathbf{K}[x]\) such that \(z \notin \mathbf{K}\) and \(\mathbf{F}' \cap \mathbf{K}[z] = \mathbf{K}\) . (Argue by contradiction, using a) and Ex. 12, b) of V, p.~149.)
\item
  Under the same hypotheses as in b), suppose that \(F' = K(y)\) , where \(y \in K[x]\) ; show that there exists \(u \in K[x]\) such that \(u \notin K\) and \(F' \cap K(u) = F' \cap K[u] = K\) . (Observe that if this were false, we should have \(F' \cap K[z] \neq K\) for all \(z \in K[x]\) such that \(z \notin K\) , using Ex. 12, c) of V, p.~149.)
\end{enumerate}

\begin{enumerate}
\def\labelenumi{\arabic{enumi})}
\setcounter{enumi}{1}
\item
  Let K be a field of characteristic p \textgreater{} 0 and f an irreducible monic polynomial of K{[}X{]}. Show that in K{[}X{]} the polynomial \(f(\mathbf{X}^{p})\) is either irreducible or the p-th power of an irreducible polynomial, depending on whether or not there exists a coefficient of f not belonging to \(K^{p}\) (decompose \(f(\mathbf{X}^{p})\) into factors of the first degree in \(\Omega[X]\) , where \(\Omega\) is an algebraically closed extension of K).
\item
  Let K be a field of characteristic p \textgreater{} 2 and let F be the field of rational fractions \(\mathrm{K}(\mathbf{X}, \mathbf{Y})\) . Let \(\mathrm{E} = \mathrm{F}(\theta)\) be an extension of F generated by a root \(\theta\) of the polynomial
\end{enumerate}

\[
f (\mathbf {Z}) = \mathbf {Z} ^ {2 p} + \mathbf {X Z} ^ {p} + \mathbf {Y}
\]

of F{[}Z{]}. Show that \([E:F(\theta^{p})]=p\) , but that E contains no p-radical element over F not in F. (Firstly remark that f is irreducible in F{[}Z{]}; if \(\beta\in E\) existed such that \(\beta^{p}\in F\) , \(\beta\notin F\) , f would be reducible in F( \(\beta\) ){[}Z{]}; using Ex. 2, show that then \(X^{1/p}\) and \(Y^{1/p}\) would belong to E and we would have \([E:F]\geqslant p^{2}\) .) (See Ex. 1, V, §7).

\BourbakiExerciseGroup

\begin{enumerate}
\def\labelenumi{\arabic{enumi})}
\item
  Let K be a field and A a commutative K-algebra. For A to be etale it is necessary and sufficient that there should exist a finite subset S of A generating the K-algebra A and such that for all \(x \in S\) the subalgebra \(K[x]\) generated by x is etale. (If S generates A, the K-algebra A is isomorphic to a quotient of the tensor product of the \(K[x]\) , where x runs over S.)
\item
  Let \(\Gamma\) be a finite commutative monoid. Consider the following conditions \(E_0\) and \(E_p\) where \(p\) denotes a prime number.
\end{enumerate}

(E0) If \(x \in \Gamma\) and if \(e, d\) are integers \(\geqslant 0\) , the relation \(x^{e + d} = x^d\) implies the relation \(x^{e + 1} = x\) .

\((E_p)\) If \(x, y \in \Gamma\) , the relation \(x^p = y^p\) implies \(x = y\) .

Let \(\mathbf{K}\) be a field of characteristic \(p (\geqslant 0)\) and let \(\mathbf{A}\) be the K-algebra \(\mathbf{K}^{(\Gamma)}\) of the monoid \(\Gamma\) . \(a)\) For \(\mathbf{A}\) to be etale it is necessary and sufficient that the condition \((\mathrm{E}_p)\) be satisfied. (Use Ex. 1 for the sufficiency. For the necessity observe for \(p = 0\) , that \(x^{e+1} - x\) is nilpotent if \(x^{e+d} = x^d\) ; for \(p > 0\) use the Cor. of V, p.~35.)

\begin{enumerate}
\def\labelenumi{\alph{enumi})}
\setcounter{enumi}{1}
\item
  Let X be the set of homomorphisms of \(\Gamma\) into the multiplicative monoid of K. We have Card \(X \leqslant\) Card \(\Gamma\) . If K is algebraically closed, then Card X = Card \(\Gamma\) holds if and only if condition \((E_{p})\) is satisfied.
\item
  If \(\Gamma\) is a group, condition \((\mathbf{E}_0)\) always holds; if \(p > 0\) , condition \((\mathbf{E}_p)\) is equivalent to saying that each element of \(\Gamma\) has order prime to \(p\) , or also that Card \(\Gamma\) is not divisible by \(p\) .
\end{enumerate}

\BourbakiExerciseGroup

\begin{enumerate}
\def\labelenumi{\arabic{enumi})}
\tightlist
\item
  \begin{enumerate}
  \def\labelenumii{\alph{enumii})}
  \tightlist
  \item
    Let \(\mathbf{K}\) be a field of characteristic \(p > 0\) . An algebraic extension \(\mathbf{F}\) of \(\mathbf{K}\) of finite degree is said to be exceptional if it is not separable over \(\mathbf{K}\) and if it contains no element \(x \notin K\) which is \(p\) -radical over \(K\) , in other words if \(F^p \cap K = K^p\) (V, p.~151, Ex. 3). Show that if \(F\) is exceptional and if \(E = F_s\) , \(z \in F\) , \(z \notin E\) and \(z^p \in E\) , then for each \(t \in K\) we have \(z^p \notin E^p(t)\) (observe that otherwise we should have \(E^p(t) = E^p(t^p)\) and so \(t \in F^p\) ).
  \end{enumerate}
\end{enumerate}

\begin{enumerate}
\def\labelenumi{\alph{enumi})}
\setcounter{enumi}{1}
\tightlist
\item
  Conversely, let E be a separable algebraic extension of K of finite degree, and let \(x \in E\) be such that for each \(t \in K\) we have \(x \notin E^p(t)\) . Let \(F = E(x^{1/p})\) , so that \(E = F_s\) . Show that F is an exceptional extension of K. (If there existed \(y \in F\) such that \(y \notin E\) and \(y^p = z \in K\) , show that we would have \(z \in E^p(x)\) ; this would imply \(z \in E^p\) , hence \(y \in E\) , which is a contradiction.)
\end{enumerate}

A separable extension E of K of finite degree is called strongly separable if the union of the fields \(E^{p}(t)\) , where t ranges over K, is distinct from E; it thus comes to the same to say that \(E = F_{s}\) , where F is a inseparable exceptional extension of K. A strongly separable extension of K cannot be a perfect field.

\begin{enumerate}
\def\labelenumi{\alph{enumi})}
\setcounter{enumi}{2}
\tightlist
\item
  Show that if E is a strongly separable extension of K then it is impossible that \([E:E^{p}]=p\) . (Argue by contradiction, showing that \([E:E^{p}]=p\) would imply \(\mathrm{E}^{p}(t)=\mathrm{E}\) for all \(t\in K\) such that \(t\notin K^{p}\) ; observe on the other hand that if \(K\subset E^{p}\) , E cannot be strongly separable over K unless it is perfect.)
\end{enumerate}

\(d\) ) Show that if \([\mathbf{K}:\mathbf{K}^p ] = p\) , then there is no strongly separable extension of \(\mathbf{K}\) , and so neither is there an exceptional extension of \(\mathbf{K}\) .

\begin{enumerate}
\def\labelenumi{\arabic{enumi})}
\setcounter{enumi}{1}
\tightlist
\item
  Let \(k\) be an imperfect field of characteristic \(p, a \in k\) an element such that \(a \notin k^p\) , \(F\) the field \(k(X)\) of rational fractions in one indeterminate over \(k\) ; put \(y = X^{p^2} / (X^p + a)\) and \(K = k(y)\) .
\end{enumerate}

\begin{enumerate}
\def\labelenumi{\alph{enumi})}
\item
  Show that \([\mathbf{F}:\mathbf{K}] = p^2\) , so that \(\mathrm{T}^{p^2} - y\mathrm{T}^p - ay\) is the minimal polynomial of \(\mathbf{X}\) in \(\mathbf{K}[\mathrm{T}]\) ; deduce that the relative separable closure of \(\mathbf{K}\) in \(\mathbf{F}\) is \(\mathrm{E} = k(\mathbf{X}^p)\) (cf.~V, p.~148, Ex. 11).
\item
  Show that F is an exceptional extension (Ex. 1) of K. (Argue by contradiction, assuming \(b \in F\) such that \(b \notin K\) and \(b^{p} \in K\) . Observe that \(\mathbf{K}(b) = k(z)\) , where \(z = r(\mathbf{X}) / s(\mathbf{X})\) with \(\deg r = p\) , \(\deg s \leqslant p\) and r, s being relatively prime polynomials in \(k[X]\) ; show that \(z^{p} \in K\) and that the minimal polynomial of X over K is \(r(T)^{p} - z^{p}s(T)^{p}\) up to a factor from K. Conclude, using a) that we would have \(a \in k \cap F^{p} = k^{p}\) contrary to the hypothesis.)
\end{enumerate}

\begin{enumerate}
\def\labelenumi{\arabic{enumi})}
\setcounter{enumi}{2}
\tightlist
\item
  Let K be a field of characteristic p \textgreater{} 0, E an extension of finite degree of K, \(E_{r}\) and \(E_{s}\) the largest p-radical extension and separable extension of K contained in E. Show that \(E_{s} \otimes_{K} E_{r}\) is isomorphic to the relative separable closure of \(E_{r}\) in E, with which it will be identified; if it is not equal to E, then E is an exceptional extension of \(E_{r}\) and \(E_{s} \otimes_{K} E_{r}\) is a strongly separable extension of \(E_{r}\) .
\end{enumerate}

Conversely, let E be an exceptional extension of a field L and let K be a subfield of L such that L is p-radical over K and \([L:K]<+\infty\) ; then L is the largest p-radical extension of K contained in E.

\begin{enumerate}
\def\labelenumi{\arabic{enumi})}
\setcounter{enumi}{3}
\tightlist
\item
  Let \(\mathbf{K}\) be a field, \(\mathbf{E} = \mathbf{K}(x)\) a separable algebraic extension of \(\mathbf{K}\) and \(f\) the minimal polynomial of \(x\) in \(\mathbf{K}[\mathbf{X}]\) .
\end{enumerate}

\begin{enumerate}
\def\labelenumi{\alph{enumi})}
\item
  Let \(\mathbf{F}\) be an extension of \(\mathbf{K}\) ; in the ring \(\mathbf{F}[\mathbf{X}]\) the polynomial \(f\) decomposes into a product \(f_{1}f_{2}\ldots f_{r}\) of irreducible and separable pairwise distinct polynomials. Show that the algebra \(\mathbf{E} \otimes_{\mathbf{K}} \mathbf{F}\) is isomorphic to the product of fields \(\mathbf{F}[\mathbf{X}] / (f_j)\) ( \(1 \leqslant j \leqslant r\) ).
\item
  Suppose that \(\mathbf{F} = \mathbf{K}(y)\) is a separable algebraic extension of \(\mathbf{K}\) ; let \(g\) be the minimal polynomial of \(y\) in \(\mathbf{K}[\mathbf{X}]\) , and let \(g = g_1g_2\ldots g_s\) be the decomposition of \(g\) into a product of irreducible polynomials in E{[}X{]}. Show that r = s and that if \(m = \deg f\) , \(n = \deg g\) , \(m_j = \deg f_j\) , \(n_j = \deg g_j\) we can, after permuting the \(g_j\) if necessary, suppose that \(m_j / n_j = m / n\) for \(1 \leq j \leq r\) .
\end{enumerate}

\begin{enumerate}
\def\labelenumi{\arabic{enumi})}
\setcounter{enumi}{4}
\tightlist
\item
  Let \(\mathbf{K}\) be a finite field with \(q\) elements, and let \(\mathbf{A}\) be the K-algebra \(\mathbf{K}^n\) . For \(\mathbf{A}\) to be generated by a single element, it is necessary and sufficient that \(n \leqslant q\) .
\end{enumerate}

\BourbakiExerciseGroup

\begin{enumerate}
\def\labelenumi{\arabic{enumi})}
\item
  Let \(\mathbf{K}\) be a commutative ring, \(f(\mathbf{X})\) a monic polynomial in \(\mathbf{K}[\mathbf{X}]\) , \(\mathbf{A}\) the K-algebra \(\mathbf{K}[\mathbf{X}] / (f)\) and \(x\) the residue class of \(\mathbf{X}\) in \(\mathbf{A}\) . For each \(a \in \mathbf{K}\) we have \(f(a) = \mathrm{N}_{\mathrm{A} / \mathrm{K}}(a - x)\) .
\item
  Let K be a field and A a K-algebra of finite degree n.~For each integer m relatively prime to n we have \(A^{*m} \cap K^{*} = K^{*m}\) ; in other words, if \(x \in K^{*}\) and if there exists an element y of A such that \(x = y^{m}\) , then there exists an element z in K such that \(x = z^{m}\) .
\end{enumerate}

\BourbakiExerciseGroup

\begin{enumerate}
\def\labelenumi{\arabic{enumi})}
\tightlist
\item
  The polynomial \(X^{2}-2\) is irreducible in Q{[}X{]} (cf.~I, p.~51, Th. 7); let \(\alpha\) be one of its roots in an algebraic closure \(\Omega\) of Q. Show that the polynomial \(X^{2}-\alpha\) is irreducible over \(E = Q(\alpha)\) ; let \(\beta\) be one of the roots of this polynomial in \(\Omega\) and let \(F = E(\beta) = Q(\beta)\) . Show that F is not a quasi-Galois extension of Q; what is the quasi-Galois extension of Q generated by F?
\end{enumerate}

* 2) Let E be the relative algebraic closure of Q in R (« the field of real algebraic numbers »); if u is an automorphism of E, then \(u(x) = x\) for all \(x \in \mathbf{Q}\) ; show that if \(y \in E\) is such that \(y > 0\) , then also \(u(y) > 0\) . Deduce that \(u(y) = y\) for all \(y \in E\) . *

\begin{enumerate}
\def\labelenumi{\arabic{enumi})}
\setcounter{enumi}{2}
\item
  Show that every algebraic extension of a field \(\mathbf{K}\) , generated by a set of elements of degree 2 over \(\mathbf{K}\) is quasi-Galois over \(\mathbf{K}\) .
\item
  Let K be a field, f an irreducible polynomial in K{[}X{]}, separable over K and of degree n and let \(\alpha_{i}\) ( \(1 \leqslant i \leqslant n\) ) be its roots in an algebraic closure \(\Omega\) of K. Let g be a polynomial in K{[}X{]}, and h an irreducible factor in K{[}X{]} of the polynomial \(f(g(\mathbf{X}))\) . Show that the degree of h is a multiple rn of n and that h has exactly r roots in common with each of the polynomials \(g(\mathbf{X}) - \alpha_{i}\) ( \(1 \leqslant i \leqslant n\) ) of \(\Omega[\mathbf{X}]\) (consider the conjugates over K of a root of h).
\item
  \begin{enumerate}
  \def\labelenumii{\alph{enumii})}
  \tightlist
  \item
    Let K be a field, \(\Omega\) an algebraic closure of K and E, F two subextensions of \(\Omega\) . Show that every composite extension of E and F (V, p.~12) is isomorphic to a composite extension of the form \((L_{w}, 1, w)\) , where w denotes a K-isomorphism of F onto a subfield of \(\Omega\) , 1 is the identity mapping of E and \(L_{w}\) is the subfield of \(\Omega\) generated by \(E \cup w(F)\) .
  \end{enumerate}
\end{enumerate}

\begin{enumerate}
\def\labelenumi{\alph{enumi})}
\setcounter{enumi}{1}
\item
  Deduce from a) that if F is of finite degree n over K, there are at most n composite extensions of E and F that are pairwise non-isomorphic.
\item
  Suppose that E is a quasi-Galois extension of K and F a subextension of E. Show that every composite extension of E and F is isomorphic to a composite extension of the form \((E, 1, v)\) , where v is a K-isomorphism of F onto a subfield of E.
\end{enumerate}

\begin{enumerate}
\def\labelenumi{\arabic{enumi})}
\setcounter{enumi}{5}
\tightlist
\item
  \begin{enumerate}
  \def\labelenumii{\alph{enumii})}
  \tightlist
  \item
    Let K be a field of characteristic exponent p, L an extension of K, \(\Omega\) an algebraic closure of L and \(\bar{K}\) the relative algebraic closure of K in \(\Omega\) . Show that \(\bar{K}\) and \(L^{p^{-\infty}}\) are linearly disjoint over \(E = \bar{K} \cap L^{p^{-\infty}}\) . (Let \((a_{\lambda})\) be a basis of the vector space \(L^{p^{-\infty}}\) over E; consider a relation \(\sum_{\lambda} c_{\lambda} a_{\lambda} = 0\) , where \(c_{\lambda} \in \bar{K}\) and the number of non-zero
  \end{enumerate}
\end{enumerate}

\(c_{\lambda}\) is as small as possible, and apply an L-automorphism of \(\Omega\) to this relation.)

\begin{enumerate}
\def\labelenumi{\alph{enumi})}
\setcounter{enumi}{1}
\tightlist
\item
  Deduce from a) that if E and F are two extensions of K such that E is algebraic over K and K is relatively algebraically closed in F, then any two composite extensions of E and F are isomorphic.
\end{enumerate}

\begin{enumerate}
\def\labelenumi{\arabic{enumi})}
\setcounter{enumi}{6}
\tightlist
\item
  Let K be a field of characteristic 0, \(\alpha\) , \(\beta\) , \(\xi\) three elements of an algebraic closure of K such that \([\mathbf{K}(\alpha):\mathbf{K}] = 2\) , \([\mathbf{K}(\beta):\mathbf{K}] = 3\) , \([\mathbf{K}(\xi):\mathbf{K}] = n > 1\) ; denote by \(\alpha'\) (resp. \(\beta'\) , \(\beta''\) ) the unique conjugate \(\neq \alpha\) of \(\alpha\) (resp. the conjugates \(\neq \beta\) of \(\beta\) ).
\end{enumerate}

\begin{enumerate}
\def\labelenumi{\alph{enumi})}
\item
  Show that \([\mathbf{K}(\alpha, \xi): \mathbf{K}] = n\) if \(\alpha \in \mathbf{K}(\xi)\) , \([\mathbf{K}(\alpha, \xi): \mathbf{K}] = 2n\) otherwise.
\item
  Show that \([\mathbf{K}(\beta,\xi):\mathbf{K}]=n\) if \(\beta\in\mathbf{K}(\xi)\) , \([\mathbf{K}(\beta,\xi):\mathbf{K}]=2n\) if \(\beta\in\mathbf{K}(\xi)\) and just one of \(\beta'\) , \(\beta''\) belongs to \(\mathbf{K}(\xi)\) and finally \([\mathbf{K}(\beta,\xi):\mathbf{K}]=3n\) if none of the three elements \(\beta\) , \(\beta'\) , \(\beta''\) belongs to \(\mathbf{K}(\xi)\) .
\item
  If \(d\) is the discriminant of the minimal polynomial of \(\beta\) , show that \(\mathbf{K}(\beta, \beta') = \mathbf{K}(\beta, \sqrt{d})\) .
\end{enumerate}

\(d\) ) Show that \(\mathbf{K}(\beta, \beta') = \mathbf{K}(\beta - \beta')\) .

\begin{enumerate}
\def\labelenumi{\alph{enumi})}
\setcounter{enumi}{4}
\item
  If \(n = 2\) and \(\alpha \notin \mathbf{K}(\xi)\) , show that \(\mathbf{K}(\alpha, \xi) = \mathbf{K}(\alpha + \xi)\) .
\item
  Show that \(\mathbf{K}(\alpha, \beta) = \mathbf{K}(\alpha + \beta)\) .
\end{enumerate}

\(g\) ) Show that \(\mathbf{K}(\alpha ,\beta) = \mathbf{K}(\alpha \beta)\) except when \(\alpha = aj\) and \(\beta^3 = b\) , where \(a,b\) are elements of \(\mathbf{K}\) and \(j^3 = 1\) .

\begin{enumerate}
\def\labelenumi{\arabic{enumi})}
\setcounter{enumi}{7}
\tightlist
\item
  Let K be an infinite field, E an algebraic extension of K of finite degree and \(\tau\) an automorphism of the vector K-space E such that for all \(x \in E\) , \(\tau(x)\) is conjugate to x over K. Show that \(\tau\) is a K-automorphism of the field E. (Let N be the quasi-Galois extension of K generated by E and let \(\Gamma\) be the group of K-automorphisms of N, \(\Delta\) the subgroup of \(\Gamma\) consisting of all E-automorphisms and \(\sigma_{j}\) ( \(1 \leqslant j \leqslant r\) ) a system of representatives of the left cosets of \(\Delta\) . If \((u_{i})_{1 \leqslant i \leqslant n}\) is a basis of the vector K-space E, consider the polynomial
\end{enumerate}

\[
\mathrm{F} \left(\mathrm{X} _ {1}, \dots , \mathrm{X} _ {n}\right) = \prod_ {j = 1} ^ {r} \left(\sum_ {i = 1} ^ {n} \left(\tau \left(u _ {i}\right) - \sigma_ {j} \left(u _ {i}\right)\right) \mathrm{X} _ {i}\right).
\]

\BourbakiExerciseGroup

\begin{enumerate}
\def\labelenumi{\arabic{enumi})}
\tightlist
\item
  \begin{enumerate}
  \def\labelenumii{\alph{enumii})}
  \tightlist
  \item
    Let K be a field, f a separable polynomial in K{[}X{]} and N a splitting field of f.~Show that for the group \(\operatorname{Gal}(N/K)\) to operate transitively on the set of roots of f in N it is necessary and sufficient for f to be irreducible in K{[}X{]}.
  \end{enumerate}
\end{enumerate}

\begin{enumerate}
\def\labelenumi{\alph{enumi})}
\setcounter{enumi}{1}
\tightlist
\item
  Suppose further that \(f\) is irreducible in \(\mathbf{K}[\mathbf{X}]\) ; let \(x\) be a root of \(f\) in \(\mathbf{N}\) . For there to exist no field \(\mathbf{E}\) such that \(\mathbf{K} \subset \mathbf{E} \subset \mathbf{K}(x)\) , distinct from \(\mathbf{K}\) and \(\mathbf{K}(x)\) it is necessary and sufficient that \(\operatorname{Gal}(\mathbf{N} / \mathbf{K})\) , qua transitive permutation group of the set of roots of \(f\) in \(\mathbf{N}\) , should be primitive (I, p.~142, Ex. 13).
\end{enumerate}

\begin{enumerate}
\def\labelenumi{\arabic{enumi})}
\setcounter{enumi}{1}
\tightlist
\item
  Let K be a field, \(f_{0}\) an irreducible and separable polynomial in K{[}X{]}, of degree n; let \(\alpha_{i}\) ( \(1 \leqslant i \leqslant n\) ) be its roots in a splitting field N of \(f_{0}\) . Let F be the field of rational fractions \(\mathrm{N}(\mathbf{X}_{1}, \mathbf{X}_{2}, \ldots, \mathbf{X}_{n})\) , E the subfield \(\mathrm{K}(\mathbf{X}_{1}, \mathbf{X}_{2}, \ldots, \mathbf{X}_{n})\) of F; F is a Galois extension of E and \(\operatorname{Gal}(\mathrm{E}/\mathrm{F})\) is canonically isomorphic to \(\operatorname{Gal}(\mathrm{N}/\mathrm{K})\) (V, p.~71, Th. 5; V, p.~154, Ex. 8). Show that if \(\theta = \alpha_{1}\mathbf{X}_{1} + \alpha_{2}\mathbf{X}_{2} + \cdots + \alpha_{n}\mathbf{X}_{n}\) then \(\mathrm{F} = \mathrm{E}(\theta)\) and the minimal polynomial \(g_{1}\) of \(\theta\) in E{[}T{]} is an irreducible factor of the polynomial
\end{enumerate}

\[
f (T) = \prod_ {\sigma \in \mathfrak {S} _ {n}} \left(T - \alpha_ {1} X _ {\sigma (1)} - \alpha_ {2} X _ {\sigma (2)} - \dots - \alpha_ {n} X _ {\sigma (n)}\right).
\]

Each permutation \(\sigma \in \mathfrak{S}_n\) defines a K-automorphism of E, again denoted by \(\sigma\) , by the condition \(\sigma(X_i) = X_{\sigma(i)}\) for \(1 \leqslant i \leqslant n\) ; show that \(\operatorname{Gal}(N/K)\) is isomorphic to the subgroup of \(\mathfrak{S}_n\) consisting of permutations \(\sigma\) such that \(\sigma(g_1) = g_1\) . Further, if \(f = g_1g_2 \ldots g_r\) is a decomposition of \(f\) into irreducible factors in E{[}T{]}, show that for every index \(j\) such that \(1 \leqslant j \leqslant r\) there exists a permutation \(\sigma_j \in \mathfrak{S}_n\) such that \(\sigma_j(g_1) = g_j\) .

\begin{enumerate}
\def\labelenumi{\arabic{enumi})}
\setcounter{enumi}{2}
\item
  Let N be a Galois extension of a field K of infinite degree over K. Show that the cardinal of the Galois group \(\operatorname{Gal}(N/K)\) is greater than that of the set of subsets of N (use V, p.~61 and 62). Deduce that there exist non-closed subgroups of \(\operatorname{Gal}(N/K)\) .
\item
  Let \(\mathbf{N}\) be a Galois extension of a field \(\mathbf{K}\) , \((\mathrm{E}_i)_{i\in I}\) a family of subextensions of \(\mathbf{N}\) , Galois over \(\mathbf{K}\) , such that: \(1^{\circ}\) for each index \(i\) , if \(\mathbf{F}_i\) is the field generated by the \(\mathbf{E}_j\) for \(j\neq i\) , then \(\mathbf{E}_i\cap \mathbf{F}_i = \mathbf{K}\) ; \(2^{\circ}\) \(\mathbf{N}\) is generated by the union of the \(\mathbf{E}_i\) .
\end{enumerate}

\begin{enumerate}
\def\labelenumi{\alph{enumi})}
\tightlist
\item
  Show that \(\mathbf{N}\) is isomorphic to the tensor product \(\otimes_{i\in I}\mathrm{E}_i\) (III, p.~470) (define an
\end{enumerate}

isomorphism of this tensor product onto N, using Th. 5 of V, p.~71).

\begin{enumerate}
\def\labelenumi{\alph{enumi})}
\setcounter{enumi}{1}
\tightlist
\item
  Show that \(\operatorname{Gal}(\mathbf{N}/\mathbf{K})\) is isomorphic to the product of the topological groups \(\operatorname{Gal}(\operatorname{E}_{i}/\operatorname{K})\) .
\end{enumerate}

\begin{enumerate}
\def\labelenumi{\arabic{enumi})}
\setcounter{enumi}{4}
\item
  Show that every compact totally disconnected group is isomorphic to the group Gal(N/K) of a Galois extension of a field (using Gen.~Top. III, p.~325, Ex. 3, reduce to the case where the group is a product of finite groups and use Ex. 4 as well as V, p.~59, Example 5).
\item
  Let N be a Galois extension of a field K, of infinite degree, and let \(\rho, \sigma\) be two distinct elements of \(\operatorname{Gal}(N/K)\) . If J is the set of \(x \in N\) such that \(\rho(x) \neq \sigma(x)\) , show that \(\mathrm{K}(\mathrm{J})\) is an extension of K of infinite degree. (Observe that a field cannot be the union of two of its subfields unless it is equal to one of them.)
\item
  Let \(\mathbf{N}\) be a Galois extension of a field \(\mathbf{K}\) such that \(\operatorname{Gal}(\mathbf{N} / \mathbf{K})\) is isomorphic to the symmetric group \(\mathfrak{S}_n\) (V, p.~59, Example 5), \(n\) being a non-prime integer; identify \(\operatorname{Gal}(\mathbf{N} / \mathbf{K})\) and \(\mathfrak{S}_n\) . Let \(\Gamma_1\) be the subgroup of \(\mathfrak{S}_n\) of order \((n - 1)!\) leaving the number 1 invariant and let \(\Gamma_2\) be the cyclic subgroup of \(\mathfrak{S}_n\) of order \(n\) generated by the cycle \((1, 2, \ldots, n)\) (I, p.~142, Ex. 12). Let \(E_1, E_2\) be the field of invariants of the subgroups \(\Gamma_1, \Gamma_2\) respectively. Show that \(E_1\) and \(E_2\) are linearly disjoint over \(\mathbf{K}\) , that there exists no field \(F\) other than \(\mathbf{K}\) and \(E_1\) such that \(K \subset F \subset E_1\) but that there exist fields \(F'\) distinct from \(E_2\) and \(N\) such that \(E_2 \subset F' \subset N\) (use Ex. 13 of I, p.~142 and Ex. 14, I, p.~143).
\item
  Let K be a field and N a Galois extension of K of finite degree n.
\end{enumerate}

\begin{enumerate}
\def\labelenumi{\alph{enumi})}
\item
  Let \(p\) be a prime number dividing \(n\) and let \(p^r\) be the highest power of \(p\) dividing \(n\) . Show that there exist subextensions \(L_i\) of \(N\) for \(0 \leqslant i \leqslant r\) such that for \(0 \leqslant i \leqslant r - 1\) , \(L_i\) contains \(L_{i+1}\) , \(L_0 = N\) , \(L_i\) is a Galois extension of degree \(p\) of \(L_{i+1}\) and \([L_r: K]\) is not divisible by \(p\) (cf.~I, p.~77, Th. 1 and I, p.~78, Th. 2).
\item
  Let \(x \in \mathbf{N}\) be such that \(\mathbf{N}\) is generated by \(x\) and its conjugates. Show that there exists a subextension \(\mathbf{E}\) of \(\mathbf{N}\) such that \(x \notin \mathbf{E}\) and \([\mathbf{N} : \mathbf{E}] = p'\) . (Consider the Sylow \(p\) -groups \(\mathrm{H}_j\) of \(\operatorname{Gal}(\mathbf{N} / \mathbf{K})\) and the (normal) subgroup \(\mathbf{H}\) generated by their union; if \(F_j\) (resp. \(\mathbf{F}\) ) is the subfield of invariants of \(\mathrm{H}_j\) (resp. \(\mathbf{H}\) ) show that \(x \notin \mathbf{F}\) and deduce that there is an index \(j\) such that \(x \notin F_j\) ; then we can take \(E = F_{j}\) .) Conclude that there exists a subextension \(L\) of \(\mathbf{N}\) such that \(L(x)\) is a Galois extension of degree \(p\) of \(L\) (use I, p.~77, Prop. 12).
\item
  Conversely, let \(\Omega\) be an algebraic closure of \(N\) and suppose that there exists a subextension \(L\) of \(\Omega\) such that \(L(x)\) is Galois of degree \(p\) over \(L\) . Show that then \(p\) divides \(n\) . (Observe that \(N \cap L(x) = (N \cap L)(x)\) and that \((N \cap L)(x)\) is a Galois extension of degree \(p\) of \(N \cap L\) .)
\end{enumerate}

\begin{enumerate}
\def\labelenumi{\arabic{enumi})}
\setcounter{enumi}{8}
\tightlist
\item
  Let K be a field of characteristic \(p\) , \(\Omega\) an algebraically closed extension of K, \(x\) an element of \(\Omega\) such that \(x \notin K\) and M a maximal element of the set of subextensions of \(\Omega\) not containing \(x\) (V, p.~148, Ex. 9), so that \(\Omega\) is an algebraic extension of M (loc. cit.). a) Show that if M is not a perfect field, \(\Omega\) is a \(p\) -radical extension of M, union of the fields \(M(x^{p^{-n}})\) for \(n \geqslant 0\) (observe that for every element \(y \in \Omega\) \(p\) -radical over M, the field \(M(y)\) contains \(M(x)\) ); deduce that we necessarily have \([M(x):M] = p\) , because every element of \(\Omega\) separable over M is necessarily in M).
\end{enumerate}

\begin{enumerate}
\def\labelenumi{\alph{enumi})}
\setcounter{enumi}{1}
\tightlist
\item
  Show that if M is perfect, \(\mathbf{M}(x)\) is a Galois extension of M, of prime degree q, and for every subextension N of \(\Omega\) which is Galois of finite degree over M, the degree \([N:M]\) is a power of q. (Take for q a prime number dividing \([\mathbf{M}(x):\mathbf{M}]\) ; note that every extension of M distinct from M and contained in \(\Omega\) contains \(\mathbf{M}(x)\) , and use Ex. 8, a).)
\end{enumerate}

\begin{enumerate}
\def\labelenumi{\arabic{enumi})}
\setcounter{enumi}{9}
\tightlist
\item
  Let \(\mathbf{K}\) be a perfect field such that for every integer \(n > 1\) there exists a Galois extension \(\mathbf{N}\) of \(\mathbf{K}\) such that \(\operatorname{Gal}(\mathbf{N} / \mathbf{K})\) is isomorphic to the alternating group \(\mathfrak{A}_n\) ; for example, for every field \(k\) of characteristic 0 the field \(\mathbf{K} = k(\mathbf{X}_n)_{n \in \mathbb{N}}\) of rational fractions in an infinity of indeterminates has this property (V, p.~59, Example 5).
\end{enumerate}

\begin{enumerate}
\def\labelenumi{\alph{enumi})}
\item
  Let A be a finite set of integers \textgreater{} 1 and let \(\Omega\) be an algebraic closure of K. An element \(x \in \Omega\) is said to be A-constructible if there exists an increasing sequence \(K = E_{0} \subset E_{1} \subset \cdots \subset E_{r} \subset \Omega\) of extensions of K such that each of the degrees \([E_{j}: E_{j-1}]\) belongs to A for \(1 \leqslant j \leqslant r\) and \(x \in E_{r}\) . Show that there exists a maximal element L in the set of subextensions \(E \subset \Omega\) such that each \(x \in E\) is A-constructible. For each \(x \notin L\) show that the degree \([L(x): L]\) belongs to N - A.
\item
  Let \(a\) be the greatest element of \(\mathbf{A}\) and let \(n\) be an integer such that \(n \geqslant \max (a + 1,5)\) . Show that there exists an extension \(\mathbf{M}\) of \(\mathbf{L}\) such that \([\mathbf{M}:\mathbf{L}] = n\) . (Observe firstly that there exists an irreducible polynomial \(f \in \mathbf{K}[\mathbf{K}]\) of degree \(n\) such that if \(\mathbf{N}\) is the splitting field of \(f\) , \(\operatorname{Gal}(\mathbf{N} / \mathbf{K})\) is isomorphic to \(\mathfrak{A}_n\) ; let \(y \in \Omega\) be a root of \(f\) in \(\Omega\) and let \(E = K(y)\) . Show that we necessarily have \(E \cap L = K\) and therefore \(L(E)\) is the desired field. Argue by contradiction, noting that the elements of \(E \cap L\) are A-constructible and that the relation \(E \cap L \neq K\) would imply the existence of a normal subgroup \(\Gamma\) of \(\mathfrak{A}_n\) such that \(1 < (\mathfrak{A}_n:\Gamma) < n\) , which would contradict the simplicity of \(\mathfrak{A}_n\) (I, p.~143, Ex. 16).
\end{enumerate}

\begin{enumerate}
\def\labelenumi{\arabic{enumi})}
\setcounter{enumi}{10}
\tightlist
\item
  \begin{enumerate}
  \def\labelenumii{\alph{enumii})}
  \tightlist
  \item
    Let E be a separable extension of a field K of finite degree and N the Galois extension of K generated by E. If \(\alpha\) is an element of N whose conjugates form a normal basis of N over K and if \(\beta = \operatorname{Tr}_{\mathbb{N}/\mathbb{E}}(\alpha)\) , show that \(\operatorname{E} = \operatorname{K}(\beta)\) .
  \end{enumerate}
\end{enumerate}

\begin{enumerate}
\def\labelenumi{\alph{enumi})}
\setcounter{enumi}{1}
\tightlist
\item
  Let N be a Galois extension of K of finite degree and let \(\alpha\) be an element of N whose conjugates form a normal basis over K. If L is a subextension of N which is Galois over K and if \(\beta = \mathrm{Tr}_{\mathbb{N}/\mathbb{L}}(\alpha)\) , show that the conjugates of \(\beta\) form a normal basis of L over K.
\item
  Let L, M be two Galois extensions of K of finite degree, such that \(L \cap M = K\) ; let \(\alpha\) (resp. \(\beta\) ) be an element of L (resp. M) whose conjugates over K form a normal basis of L (resp. M) over K; show that the conjugates over K of \(\alpha\beta\) form a normal basis of \(\mathrm{K}(\mathrm{L} \cup \mathrm{M})\) .
\end{enumerate}

\begin{enumerate}
\def\labelenumi{\arabic{enumi})}
\setcounter{enumi}{11}
\tightlist
\item
  Let K be an infinite field, N a Galois extension of degree n of K, \(f \in K[X]\) the minimal polynomial of an element \(\theta\) of N such that \(\mathbf{N} = \mathbf{K}(\theta)\) . The group \(\Gamma = \operatorname{Gal}(N/K)\) is allowed to operate on the polynomials in N{[}X{]} by operating on the coefficients of these polynomials.
\end{enumerate}

\begin{enumerate}
\def\labelenumi{\alph{enumi})}
\tightlist
\item
  Let \(g \in \mathbb{N}[\mathbf{X}]\) be the polynomial of degree \(\leqslant n - 1\) such that \(g(\theta) = 1\) , \(g(\sigma(\theta)) = 0\) for all \(\sigma \neq 1\) in \(\Gamma\) . For every subgroup \(\Delta\) of \(\Gamma\) let \(p_{\Delta} \in \mathbb{N}[\mathbf{X}]\) be the polynomial \(\det(\sigma\tau(g))_{\sigma \in \Delta, \tau \in \Delta}\) . Show that
\end{enumerate}

\[
p _ {\Delta} ^ {2} \equiv \sum_ {\sigma \in \Delta} \sigma (g) \mod f
\]

and deduce that \(p_{\Delta} \neq 0\) .

\begin{enumerate}
\def\labelenumi{\alph{enumi})}
\setcounter{enumi}{1}
\tightlist
\item
  Deduce that there exists an element \(\alpha \in K\) such that for every subgroup \(\Delta\) of \(\Gamma\) the conjugates of \(g(\alpha)\) over the subfield \(k(\Delta)\) form a normal basis of N over \(k(\Delta)\) .
\end{enumerate}

\begin{enumerate}
\def\labelenumi{\arabic{enumi})}
\setcounter{enumi}{12}
\item
  Let E be an algebraic extension of a field K, \(E_{s}\) the largest separable extension and \(E_{r}\) the largest p-radical extension of K contained in E. In an algebraic closure of K let N be the quasi-Galois extension of K generated by E; then the largest separable extension of K contained in N is the quasi-Galois extension \(N_{s}\) generated by \(E_{s}\) . For \(E_{r}\) to be the largest p-radical extension of K contained in N, it is necessary and sufficient that E should be separable over \(E_{r}\) (cf.~V, p.~152, Ex. 3).
\item
  Let \(\mathbf{E}\) be an algebraic extension of a field \(\mathbf{K}\) , \(\Gamma\) the group of \(\mathbf{K}\) -automorphisms of \(\mathbf{E}\) and \(\mathbf{S} \subset \mathbf{E}\) the field of invariants of \(\Gamma\) .
\end{enumerate}

\begin{enumerate}
\def\labelenumi{\alph{enumi})}
\item
  Show that for E to be quasi-Galois over K it is necessary and sufficient for S to be a p-radical extension of K.
\item
  Let \(S_{s}\) be the largest separable extension of K contained in S. Show that \(S_{s}\) is the smallest of the fields F such that \(K \subset F \subset E\) and E is a quasi-Galois extension of F.
\item
  Let \(E_{s}\) be the largest separable extension of K contained in E; show that \(\mathrm{E} = \mathrm{S}(\mathrm{E}_{s})\) (observe that no K-automorphism of E other than the identity leaves all the elements of \(\mathrm{S}(\mathrm{E}_{s})\) invariant).
\end{enumerate}

\begin{enumerate}
\def\labelenumi{\arabic{enumi})}
\setcounter{enumi}{14}
\item
  Let E be a Galois extension of finite degree n of a field K whose Galois group G is nilpotent. Let P be the set of prime numbers dividing n.~Then there exist Galois subextensions \(E_{p}\) of E ( \(p \in P\) ) such that \(\operatorname{Gal}(E_{p}/K)\) is a p-group for each \(p \in P\) , and such that the canonical homomorphism \(\otimes E_{p} \to E\) is an isomorphism.
\item
  Let \(a\) be an integer and let \(x_1, x_2, x_3, x_4\) be the roots of \(P(X) = X^4 - aX - 1\) in an algebraic closure of \(\mathbf{Q}\) . We identify the Galois group \(G\) of \(P\) with a group of permutations of \(\{x_1, \ldots, x_4\}\) . Suppose that \(Q(x_1)\) contains no quadratic subextension (cf.~Ex. 1 of \(V\) , p.~154).
\end{enumerate}

\begin{enumerate}
\def\labelenumi{\alph{enumi})}
\item
  Show that \(G = S_{4}\) or \(A_{4}\) . (Use Ex. 1 of V, p.~154: P is irreducible, hence G acts transitively; further \(\mathbf{Q}(x_{1})\) contains no quadratic subextensions, hence the stabilizer H of \(x_{1}\) is not contained in a subgroup of index 2 of G, whereas \([G:H] = 4\) ; deduce that G is not a 2-group, hence its order is divisible by 3.)
\item
  Let \(\delta = \prod_{1\leqslant i < j\leqslant 4}(x_i - x_j)\) and \(d = \delta^2\) . Show that \(d = 4^4 -3^3 a^4\) ; further \(s(\delta) = \varepsilon_s\delta\) for
\end{enumerate}

all \(s \in G\) . Deduce that \(G = \mathfrak{A}_4\) if \(\delta \in Q\) , and \(G = \mathfrak{S}_4\) otherwise. Show that \(\delta\) is not rational. (We have thus constructed an infinite family of equations of degree 4 over \(Q\) with Galois group \(\mathfrak{S}_4\) .)

* 17) Let \(f\) be an irreducible polynomial in \(\mathbf{Q}[\mathbf{X}]\) of prime degree \(p\) , having exactly two roots in \(\mathbf{C} - \mathbf{R}\) . Then the Galois group \(G\) of \(f\) is isomorphic to \(\mathfrak{S}_p\) . (G contains an element of order \(p\) and a transposition.) *

\begin{enumerate}
\def\labelenumi{\arabic{enumi})}
\setcounter{enumi}{17}
\tightlist
\item
  Let K be a field and F the field of rational fractions \(\mathrm{K}(\mathrm{T})\) . The group \(\mathrm{GL}_{2}(\mathrm{K})\) operates on F by the law \((\gamma f)(\mathrm{T}) = f\left(\frac{a\mathrm{T} + b}{c\mathrm{T} + d}\right)\) , for \(\gamma = \begin{pmatrix} a & b \\ c & d \end{pmatrix} \in \mathrm{GL}_{2}(\mathrm{K})\) and \(f \in F\) . The centre \(K^{*}\) . I of \(\mathrm{GL}_{2}(\mathrm{K})\) operates trivially, whence we obtain, by passage to quotients, a homomorphism of \(\mathrm{PGL}_{2}(\mathrm{K}) = \mathrm{GL}_{2}(\mathrm{K}) / \mathrm{K}^{*}\) . I into the automorphism group of the extension F of K. This homomorphism is bijective (V, p.~148, Ex. 11).
\end{enumerate}

\begin{enumerate}
\def\labelenumi{\alph{enumi})}
\tightlist
\item
  Let H be a subgroup of \(\mathrm{PGL}_{2}(\mathbf{K})\) of finite order h, and let E be the field of invariants of H, so that \(F = E(T)\) is a Galois extension of E with Galois group H (V, p.~66, Th. 3). Let
\end{enumerate}

\[
\mathrm{H} (\mathrm{X}) = \prod_ {s \in \mathrm{H}} (\mathrm{X} - s (\mathrm{T})) = \sum_ {i = 0} ^ {h} (- 1) ^ {i} \mathrm{S} _ {i} (\mathrm{T}) \mathrm{X} ^ {h - i}
\]

be the minimal polynomial of T over E. For each i such that \(0 \leqslant i \leqslant h\) , either \(S_{i}(T) \in K\) or \(E = K(S_{i}(T))\) . (Use Ex. 11, a) of V, p.~157.)

\begin{enumerate}
\def\labelenumi{\alph{enumi})}
\setcounter{enumi}{1}
\item
  For each \(s \in H\) let \(\begin{pmatrix} a_{s} & b_{s} \\ c_{s} & d_{s} \end{pmatrix}\) be a representative in \(GL_{2}(K)\) and let \(H_{0}\) be the subgroup of H consisting of those s for which \(c_{s} = 0\) ; let \(h_{0} = CardH_{0}\) and put \(R^{H}(T) = S_{h_{0}}(T)\) . Show that \(E = K(R^{H}(T))\) and that \(R^{H}(T) = P(T)/Q(T)^{h_{0}}\) where P is a polynomial of degree h and \(Q(T) = \prod_{s} (c_{s}T + d_{s})\) where s runs over a system of representatives of the cosets \(H_{0}s\) different from \(H_{0}\) ; the degree of Q is thus \((h/h_{0}) - 1\) .
\item
  Suppose that \(K^{*}\) contains a subgroup \(\mu_{n}\) of finite order n.~Let \(C_{n}\) be the image in \(PGL_{2}(K)\) of the subgroup of \(GL_{2}(K)\) consisting of the matrices \(\begin{pmatrix} a & 0 \\ 0 & 1 \end{pmatrix}, a \in \mu_{n}\) . Then \(C_{n}(X) = X^{n} - T^{n}\) and \(R^{C_{n}}(T) = (-1)^{n+1}T^{n}\) .
\item
  Let \(w\) be the element of \(\mathrm{PGL}_2(\mathbf{K})\) represented by \(\begin{pmatrix} 0 & 1 \\ 1 & 0 \end{pmatrix}\) . Then \(w^2 = 1\) and \(wsw^{-1} = s^{-1}\) if \(s\) is represented by a diagonal matrix. The subgroup \(\mathrm{D}_n\) of \(\mathrm{PGL}_2(\mathbf{K})\) generated by \(\mathrm{C}_n\) and \(w\) is therefore of order \(2n\) . We have \(\mathrm{D}_n(\mathbf{X}) = (\mathbf{X}^n - \mathrm{T}^n)\left(\mathbf{X}^n - \left(\frac{1}{\mathrm{T}}\right)^n\right) = \mathbf{X}^{2n} + (-1)\mathbf{R}^{\mathrm{D}_n(\mathrm{T})} + 1\) where \(\mathbf{R}^{\mathrm{D}_n(\mathrm{T})} = -((-T)^n + (-T)^{-n})\) .
\item
  Let S be the group of order 3 in \(\mathrm{PGL}_2(\mathbf{K})\) generated by the image of the matrix \(\left( \begin{array}{cc}0 & -1\\ 1 & 1 \end{array} \right)\) . We have \(S(X) = (X - T)\left(X - \frac{1}{1 - T}\right)\left(X - \frac{T - 1}{T}\right)\) and \(R^s (T) = \frac{T^3 - 3T + 1}{T^2 - 1}\) .
\item
  Suppose that K is a finite field with q elements. Let U be the image in \(\mathrm{PGL}_{2}(\mathbf{K})\) of the group of all matrices \(\begin{pmatrix}1 & 0 \\ b & 1\end{pmatrix}\) , \(b \in K\) . Then
\end{enumerate}

\[
\mathrm{U} (\mathrm{X}) = (\mathrm{X} - \mathrm{T}) ^ {q} - (\mathrm{X} - \mathrm{T}) = \mathrm{X} ^ {q} - \mathrm{X} - (\mathrm{T} ^ {q} - \mathrm{T}) \quad \text { and } \quad \mathrm{R} ^ {\mathrm{U}} (\mathrm{T}) = \mathrm{T} ^ {q} - \mathrm{T}.
\]

Let B be the group generated by \(C_{q-1}\) and U; it is of order \(q(q-1)\) . We have

\[
\mathrm{B} (\mathrm{X}) = \left(\mathrm{X} ^ {q} - \mathrm{X}\right) ^ {q - 1} - \left(\mathrm{T} ^ {q} - \mathrm{T}\right) ^ {q - 1} \quad \text { and } \quad \mathrm{R} ^ {\mathrm{H}} (\mathrm{T}) = \left(\mathrm{T} ^ {q} - \mathrm{T}\right) ^ {q - 1}.
\]

\begin{enumerate}
\def\labelenumi{\arabic{enumi})}
\setcounter{enumi}{18}
\item
  Let E be a Galois extension of a field K with Galois group G. For each \(x \in E\) denote by Gx the set of conjugates of x in E; this is a finite G-set. Let x, y be two elements of E. For the G-sets Gx and Gy to be isomorphic (that is, for a bijection \(Gx \rightarrow Gy\) to exist, compatible with the operation of G on the two sets), it is necessary and sufficient that the extensions \(\mathrm{K}(x)\) and \(\mathrm{K}(y)\) should be conjugate.
\item
  Let E be an algebraic extension of a field K such that every non-constant polynomial in K{[}X{]} has at least one root in E. Show that E is algebraically closed. (In an algebraic closure of E, hence of K, let F be an extension of finite degree of K. To show that \(F \subset E\) reduce to the case where F is quasi-Galois and then, by means of Prop. 13 of V, p.~76, to the case where F is either Galois or p-radical.)
\item
  Let \(f(\mathbf{X}) = \mathbf{X}^{n} + a_{1}\mathbf{X}^{n-1} + \cdots + a_{n}\) be a polynomial of Z{[}X{]} and p a prime number. Suppose that \(a_{i} \equiv 0 \pmod{p}\) , \(1 \leqslant i \leqslant n\) and that \(a_{n} \not\equiv 0 \pmod{p^{2}}\) . Show that \(f(\mathbf{X})\) is irreducible in Q{[}X{]}. (Let \(f(\mathbf{X}) = \mathrm{P}(\mathbf{X}) \mathrm{Q}(\mathbf{X})\) be a decomposition in Q{[}X{]}. Show that if P and Q are monic, they have integer coefficients. Reduce mod p.)
\end{enumerate}

* 22) Let \(m\) be an even integer \(>0\) , \(r\) an integer \(\geqslant 3\) and \(n_1 < \cdots < n_{r-2}\) , \(r-2\) even integers. Put \(g(\mathbf{X}) = (\mathbf{X}^2 + m)(\mathbf{X} - n_1) \ldots (\mathbf{X} - n_{r-2})\) .

\begin{enumerate}
\def\labelenumi{\alph{enumi})}
\item
  Put \(f(\mathbf{X}) = g(\mathbf{X}) - 2\) . Show that \(f(\mathbf{X})\) has at least \(r - 2\) real roots. Calculate the sum of the squares of the roots of \(f\) in \(\mathbf{C}\) . Deduce that for \(m \geqslant \sum n_i^2 / 2\) , \(f(\mathbf{X})\) is a polynomial with integer coefficients having \(r - 2\) real roots and two non-real conjugate roots.
\item
  Show that \(f(\mathbf{X})\) is irreducible in \(\mathbf{Q}[\mathbf{X}]\) . (Write \(f(\mathbf{X}) = \mathbf{X}' + \sum_{k=1}^{r} a_k \mathbf{X}'^{-k}\) ; show that
\end{enumerate}

\(a_{k}\) is even and that \(a_{r}\) is not divisible by 4; apply Ex. 21.)

\begin{enumerate}
\def\labelenumi{\alph{enumi})}
\setcounter{enumi}{2}
\item
  Show that when \(r\) is prime and \(m \geqslant \sum n_i^2 / 2\) , then the Galois group \(G\) of the splitting field of \(f\) is isomorphic to \(\mathfrak{S}_r\) (cf.~V, p.~158, Ex. 17).
\item
  Show that when \(r = 4\) and \(m \geqslant \sum n_i^2 / 2\) , the degree over \(\mathbf{Q}\) of the splitting field is equal to 8 or 24, and that it is equal to 8 if and only if \(n_1 = -n_2\) or also if and only if \(f(\mathbf{X})\) is of the form \(\mathrm{P}(\mathbf{X}^2)\) , where \(\mathbf{P}\) is a polynomial of degree 2. *
\end{enumerate}

\begin{enumerate}
\def\labelenumi{\arabic{enumi})}
\setcounter{enumi}{22}
\tightlist
\item
  Let \(k\) be a field of characteristic 2.
\end{enumerate}

\begin{enumerate}
\def\labelenumi{\alph{enumi})}
\tightlist
\item
  Let \(A_{1}, \ldots, A_{n}\) be indeterminates, K the field \(k(\mathbf{A}_{1}, \ldots, \mathbf{A}_{n})\) , E a splitting field of the polynomial \(P = X^{n} + A_{1}X^{n-1} + \cdots + A_{n} \in K[X]\) and \(X_{1}, \ldots, X_{n}\) the roots of P in E, so that \(E = k(X_{1}, \ldots, X_{n})\) (V, p.~21). We identify the Galois group of E/K with the symmetric group \(S_{n}\) acting by permutations of the variables \(X_{1}, \ldots, X_{n}\) (V, p.~59). Put
\end{enumerate}

\[
\beta \left(\mathrm{X} _ {1}, \dots , \mathrm{X} _ {n}\right) = \sum_ {i <   j} \mathrm{X} _ {i} / \left(\mathrm{X} _ {i} + \mathrm{X} _ {j}\right) \in \mathrm{E} \text { and } \mathrm{C} \left(\mathrm{A} _ {1}, \dots , \mathrm{A} _ {n}\right) = \sum_ {i <   j} \left(\mathrm{X} _ {i} \mathrm{X} _ {j}\right) / \left(\mathrm{X} _ {i} ^ {2} + \mathrm{X} _ {j} ^ {2}\right) \in \mathrm{K}.
\]

Then \(\beta \notin K\) and

\[
\boldsymbol {\beta} ^ {2} + \boldsymbol {\beta} + \mathbf {C} = 0.
\]

We have \(\sigma(\beta) = \beta\) for \(\sigma \in \mathfrak{A}_n\) and \(\sigma(\beta) = \beta + 1\) for \(\sigma \in \mathfrak{S}_n - \mathfrak{A}_n\) . Deduce that \(K(\beta)\) is the unique quadratic subextension of \(E/K\) .

\begin{enumerate}
\def\labelenumi{\alph{enumi})}
\setcounter{enumi}{1}
\tightlist
\item
  Let \(D \in Z[T_{1}, \ldots, T_{n}]\) be the discriminant of the polynomial \(X^{n} + T_{1}X^{n-1} + \cdots + T_{n}\) . Show that there exist Q, R, \(S \in Z[T_{1}, \ldots, T_{n}]\) with
\end{enumerate}

\[
\mathbf {D} = \mathbf {Q} ^ {2} - 4 \mathbf {R} + 8 \mathbf {S}
\]

and

\[
\mathrm{C} = \mathrm{C} (\mathrm{A} _ {1}, \dots , \mathrm{A} _ {n}) = \frac {\mathrm{R} (\mathrm{A} _ {1} , \dots , \mathrm{A} _ {n})}{\mathrm{Q} ^ {2} (\mathrm{A} _ {1} , \dots , \mathrm{A} _ {n})} = \frac {\mathrm{R} (\mathrm{A} _ {1} , \dots , \mathrm{A} _ {n})}{\mathrm{D} (\mathrm{A} _ {1} , \dots , \mathrm{A} _ {n})}.
\]

For every triple of elements U, V, W of \(\mathbf{Z}[T_1, \ldots, T_n]\) such that \(D = U^2 - 4V + 8W\) there exists \(M \in \mathbf{Z}[A_1, \ldots, A_n]\) such that

\[
C = \frac {V (A _ {1} , \dots , A _ {n})}{D (A _ {1} , \dots , A _ {n})} + \left(\frac {M}{D (A _ {1} , \dots , A _ {n})}\right) ^ {2} + \frac {M}{D (A _ {1} , \dots , A _ {n})}.
\]

\begin{enumerate}
\def\labelenumi{\alph{enumi})}
\setcounter{enumi}{2}
\tightlist
\item
  Let \(f = X^{n} + a_{1}X^{n-1} + \cdots + a_{n}\) be a separable polynomial of k{[}X{]} and G the Galois group of a splitting extension. Show that \((a_{1}, \ldots, a_{n})\) is substitutable in C and that the following conditions are equivalent
\end{enumerate}

\begin{enumerate}
\def\labelenumi{(\roman{enumi})}
\item
  every element of \(G\) induces an even permutation of the roots of \(f\) ;
\item
  there exists \(\alpha \in k\) such that \(C(a_{1},\ldots ,a_{n}) = \alpha^{2} + \alpha\) .
\end{enumerate}

* \(d\) ) Henceforth assume that \(k\) is finite and choose the elements U, V, W of \(\mathbf{Z}[T_1, \ldots, T_n]\) arbitrarily to satisfy the condition stated in \(b\) ).

Show that (i) and (ii) of \(c\) ) are equivalent to the following conditions:

\begin{enumerate}
\def\labelenumi{(\roman{enumi})}
\setcounter{enumi}{2}
\tightlist
\item
  \[
\operatorname{Tr} _ {k / \mathbf {F} _ {2}} (\mathbf {C}) = 0;
\]
\item
  \[
\mathrm{Tr} _ {k / \mathbf {F} _ {2}} \left(\frac {\mathbf {V} (a _ {1} , \dots , a _ {n})}{\mathbf {D} (a _ {1} , \dots , a _ {n})}\right) = 0.
\]
\end{enumerate}

\begin{enumerate}
\def\labelenumi{(\alph{enumi})}
\setcounter{enumi}{21}
\tightlist
\item
  The number of irreducible factors of \(f\) is congruent to \(n\) modulo 2 (use Ex. 22 of V, p.~165). *
\end{enumerate}

\BourbakiExerciseGroup

\begin{enumerate}
\def\labelenumi{\arabic{enumi})}
\item
  Let \(\mathbf{K}\) be a field of characteristic \(p\) , \(n\) an integer not a multiple of \(p\) , \(\mathbf{R}_n(\mathbf{K})\) the field of \(n\) -th roots of unity and \(\mathbf{E}\) a cyclic extension of \(\mathbf{R}_n(\mathbf{K})\) of degree \(n\) over \(\mathbf{K}\) ; \(\mathbf{E}\) is thus the root field of a polynomial \(X^{n}-\alpha\in\mathbb{R}_{n}(\mathbf{K})[\mathbf{X}]\) (V, p.~89, Th. 4). For E to be a Galois extension of K it is necessary and sufficient that for every automorphism \(\tau\in\operatorname{Gal}\left(\mathbb{R}_{n}(\mathbf{K})/\mathbf{K}\right)\) there exists an integer r\textgreater0 and an element \(\gamma\in\mathbb{R}_{n}(\mathbf{K})\) such that \(\tau(\alpha)=\gamma^{n}\alpha^{r}\) .
\item
  \begin{enumerate}
  \def\labelenumii{\alph{enumii})}
  \tightlist
  \item
    Let E be a splitting extension of the polynomial \(X^{3}-2\in Q[X]\) . The field E is a cyclic extension of degree 3 of \(\mathrm{R}_{3}(\mathbf{Q})\) but contains no cyclic extension of degree 3 of Q. b) The field \(\mathrm{E}=\mathrm{R}_{9}(\mathbf{Q})\) is a cyclic extension of degree 3 of \(\mathrm{R}_{3}(\mathbf{Q})\) and contains a cyclic extension of degree 3 of Q.
  \end{enumerate}
\end{enumerate}

\begin{enumerate}
\def\labelenumi{\alph{enumi})}
\setcounter{enumi}{2}
\tightlist
\item
  With the hypothesis and notation as in V, p.~154, Ex. 7, show that for \(\mathrm{K}(\alpha,\beta)\) to be Galois over K it is necessary and sufficient for d to be a square in \(\mathrm{K}(\alpha)\) ; for \(\mathrm{K}(\beta)\) to be cyclic over K it is necessary and sufficient for d to be a square in K.
\end{enumerate}

* 3) Show that for every integer \(n > 1\) there exists a cyclic extension of degree \(n\) of the field \(\mathbf{Q}\) . (Reduce to the case where \(n\) is a prime number \(q\) and use the structure of the multiplicative group \((\mathbf{Z} / q^m\mathbf{Z})^*\) (VII, p.~13). *

\begin{enumerate}
\def\labelenumi{\arabic{enumi})}
\setcounter{enumi}{3}
\tightlist
\item
  Let \(\mathbf{K}\) be a field of characteristic \(p, q\) a prime number \(\neq p\) such that \(\mathbf{K}\) contains the \(q\) -th roots of unity and let \(\zeta\) be a primitive \(q\) -th root of unity.
\end{enumerate}

\begin{enumerate}
\def\labelenumi{\alph{enumi})}
\item
  Let E be a cyclic extension of K of degree \(q^e\) and \(\sigma\) a K-automorphism of E generating Gal(E/K). Let F be the field of degree \(q^{e-1}\) over K such that \(K \subset F \subset E\) ; there exists \(\theta \in E\) , a root of an irreducible polynomial \(X^q - \alpha\) in F{[}X{]} such that \(E = F(\theta)\) and \(\theta^{\sigma^m} = \zeta\theta\) (where \(m = q^{e-1}\) ). Show that also \(E = K(\theta)\) and \(\theta^\sigma = \beta\theta\) where \(\beta \in F\) is such that \(\alpha^{\sigma-1} = \beta^q\) and \(N_{F/K}(\beta) = \zeta\) (observe that \(\theta^\sigma = \beta\theta^k\) with \(0 < k < q\) and \(\beta \in F\) , by Ex. 1; by calculating \(\theta^{\sigma^{mq}}\) , show that \(k^{mq} - 1 \equiv 0 \pmod{q}\) ) and deduce that \(k = 1\) ).
\item
  Conversely, let F be a cyclic extension of degree \(q^{e-1}\) of K with \(e \geqslant 2\) and let \(\sigma\) be a K-automorphism of F generating Gal(F/K). If there exists an element \(\beta \in F\) such that \(N_{F/K}(\beta) = \zeta\) and if \(\alpha \in F\) is such that \(\alpha^{\sigma-1} = \beta^q\) , show that for every \(c \in K^*\) the polynomial \(X^q - c\alpha\) is irreducible in F{[}X{]}. If \(\theta\) is one of the roots of this polynomial in an algebraic closure \(\Omega\) of F, show that there exists a K-homomorphism \(\bar{\sigma}\) of \(E = F(\theta)\) into \(\Omega\) , extending \(\sigma\) and such that \(\theta^{\sigma} = \beta\theta\) ; deduce that E is a cyclic extension of K of degree \(q^e\) , that \(\bar{\sigma}\) generates Gal(E/K) and that \(E = K(\theta)\) . Finally show that every cyclic extension of degree \(q^e\) of K, containing F, is the root field of a polynomial \(X^q - c\alpha\) in F{[}X{]}, for an appropriate \(c \in K^*\) (use Th. 3 of V, p.~85).
\item
  Take K to be the field Q of rational numbers. The polynomial \(X^{2} + 1\) is irreducible in Q{[}X{]}; if i is one of its roots, then \(F = \mathbf{Q}(i)\) is a cyclic extension of Q of degree 2 but there exists no cyclic extension of degree 4 of Q containing F.
\end{enumerate}

\begin{enumerate}
\def\labelenumi{\arabic{enumi})}
\setcounter{enumi}{4}
\item
  Let K be a field of characteristic p and n an integer not a multiple of p.~For a polynomial of K{[}X{]} of the form \(X^{n} - a\) to be irreducible it is necessary and sufficient that for every prime factor q of n, a should not equal the q-th power of an element of K and further when \(n \equiv 0 \pmod{4}\) , a should not be of the form \(-4c^{4}\) with \(c \in K\) . (To prove the sufficiency of the condition, reduce to the case where \(n = q^{e}\) with q prime and use Ex. 4 of V, p.~153. Argue by induction an e, determining, with the help of Ex. 4 of V, p.~153, the form of the constant term of each irreducible factor of \(X^{q^{e}} - a\) .)
\item
  Let \(\mathbf{K}\) be a field of characteristic \(p, n\) an integer not a multiple of \(p\) and \(\mathbf{N}\) a splitting field of the polynomial \(\mathbf{X}^n - a\) of \(\mathbf{K}[\mathbf{X}]\) .
\end{enumerate}

\begin{enumerate}
\def\labelenumi{\alph{enumi})}
\tightlist
\item
  Show that the group \(\operatorname{Gal}(\mathbf{N} / \mathbf{K})\) is isomorphic to a subgroup of the group \(\Gamma\) of matrices \(\left( \begin{array}{cc}x & y\\ 0 & 1 \end{array} \right)\) , where \(x\in (\mathbf{Z} / n\mathbf{Z})^*\) and \(y\in \mathbf{Z} / n\mathbf{Z}\) . If \(n\) is a prime number, the group
\end{enumerate}

Gal(N/K) is commutative only if \(\mathrm{R}_{n}(\mathrm{K})=\mathrm{K}\) or if a is the n-th power of an element of K. b) If K=Q and if n is a prime number, show that if a is not the n-th power of a rational number, then Gal(N/K) is isomorphic to the whole of \(\Gamma\) .

\begin{enumerate}
\def\labelenumi{\alph{enumi})}
\setcounter{enumi}{2}
\tightlist
\item
  If \(\mathbf{K} = \mathbf{Q}\) , determine \(\operatorname{Gal}(\mathbf{N} / \mathbf{K})\) when \(a\) is a square-free integer.
\end{enumerate}

\begin{enumerate}
\def\labelenumi{\arabic{enumi})}
\setcounter{enumi}{6}
\tightlist
\item
  An extension E of a field K is called a Kummer extension if there exists an integer n \textgreater{} 0 such that K contains a primitive n-th root of unity and E is an abelian extension with Galois group annihilated by n.
\end{enumerate}

\begin{enumerate}
\def\labelenumi{\alph{enumi})}
\item
  Let K be a field and N a Galois extension of degree m of K with Galois group \(\Gamma\) . For N to be a Kummer extension of K it is necessary and sufficient that there should exist a basis \((\theta_{\sigma})_{\sigma\in\Gamma}\) of N over K such that the elements \(\gamma_{\sigma\tau}=\theta_{\sigma}^{1-\tau}\) belong to K for each pair \((\sigma,\tau)\) . (For the necessity of the condition reduce to the case where N is cyclic of prime power degree. For the sufficiency observe that if \(\tau\in\Gamma\) is of order n, then there exists \(\sigma\in\Gamma\) such that \(\gamma_{\sigma\tau}\) is a primitive n-th root of unity.)
\item
  If \((\theta_{\sigma})_{\sigma \in \Gamma}\) is a basis satisfying condition \(a)\) , show that for an element \(x = \sum_{\sigma} a_{\sigma} \theta_{\sigma}\) of \(N\) to
\end{enumerate}

be such that the conjugates of \(x\) form a normal basis of \(\mathbf{N}\) over \(\mathbf{K}\) , it is necessary and sufficient that \(a_{\sigma} \neq 0\) for all \(\sigma \in \Gamma\) .

* 8) Let N be a Kummer extension of a field K.

\begin{enumerate}
\def\labelenumi{\alph{enumi})}
\item
  Let E be a subextension of N such that \([E:K]=q\) is a prime number. Show that there exists a decomposition \(N=N_{1}\otimes N_{2}\otimes\cdots\otimes N_{r}\) of N into cyclic extensions of K such that \(E\subset N_{j}\) for an index j.
\item
  Deduce from a) that if \(x \in N\) is such that the conjugates of x over K form a normal basis of N over K, then for each field E such that \(K \subset E \subset N\) , the conjugates of x over E form a normal basis of N over E (reduce to the case where \([E : K]\) is prime and use Ex. 7, b)). *
\end{enumerate}

\begin{enumerate}
\def\labelenumi{\arabic{enumi})}
\setcounter{enumi}{8}
\tightlist
\item
  Let \(\mathbf{K}\) be a field of characteristic \(p > 0\) .
\end{enumerate}

\begin{enumerate}
\def\labelenumi{\alph{enumi})}
\item
  Let E be a cyclic extension of degree \(p^e\) of K and let \(\sigma\) be a K-automorphism generating the group \(\operatorname{Gal}(\mathrm{E} / \mathrm{K})\) . Let F be the field of degree \(p^{e - 1}\) over K such that \(\mathbf{K} \subset \mathbf{F} \subset \mathbf{E}\) ; if \(m = p^{e - 1}\) , show that there exists \(\theta \in \mathbf{E}\) , root of an irreducible polynomial \(\mathbf{X}^p - \mathbf{X} - \alpha\) of \(\mathrm{F}[\mathbf{X}]\) such that \(\mathrm{E} = \mathrm{F}(\theta)\) and \(\sigma^m (\theta) = \theta + 1\) . Show that we also have \(\mathrm{E} = \mathrm{K}(\theta)\) and \(\sigma (\theta) = \theta + \beta\) where \(\beta \in \mathrm{F}\) is such that \(\sigma (\alpha) - \alpha = \beta^p - \beta\) and \(\mathrm{Tr}_{\mathrm{F} / \mathrm{K}}(\beta) = 1\) .
\item
  Conversely, let F be a cyclic extension of degree \(p^{e-1}\) of K (with e \textgreater{} 1), \(\sigma\) a K-automorphism of F generating Gal(F/K). Show that there exist two elements \(\alpha\) , \(\beta\) of F such that \(\operatorname{Tr}_{\mathrm{F}/\mathrm{K}}(\beta) = 1\) and \(\sigma(\alpha) - \alpha = \beta^{p} - \beta\) (use Th. 3 of V, p.~85 and the Cor. of Prop. 1 of V, p.~49). Deduce that for all \(c \in K\) the polynomial \(X^{p} - X - \alpha - c\) is irreducible in F{[}X{]}; if \(\theta\) is a root of this polynomial in an algebraic closure \(\Omega\) of F, show that there exists a K-homomorphism \(\bar{\sigma}\) of E = F( \(\theta\) ) in \(\Omega\) extending \(\sigma\) and such that \(\bar{\sigma}(\theta) = \theta + \beta\) ; conclude that E is a cyclic extension of degree \(p^{e}\) of K, that \(\bar{\sigma}\) generates Gal(E/K) and that E = K( \(\theta\) ). Finally show that every cyclic extension of degree \(p^{e}\) of K containing F is the splitting field of a polynomial \(X^{p} - X - \alpha - c\) of F{[}X{]} for a suitable \(c \in K\) .
\end{enumerate}

\begin{enumerate}
\def\labelenumi{\arabic{enumi})}
\setcounter{enumi}{9}
\tightlist
\item
  \begin{enumerate}
  \def\labelenumii{\alph{enumii})}
  \tightlist
  \item
    Let \(\mathbf{K}\) be a field whose algebraic closure is an extension of prime degree \(q\) of \(\mathbf{K}\) . Such that \(\mathbf{K}\) is perfect and that \(q\) must be distinct from the characteristic of \(\mathbf{K}\) (use Ex. 9).
  \end{enumerate}
\end{enumerate}

\begin{enumerate}
\def\labelenumi{\alph{enumi})}
\setcounter{enumi}{1}
\tightlist
\item
  Show that K contains the q-th roots of unity and that \(\Omega\) is the splitting field of an irreducible polynomial \(X^{q}-a\) of K{[}X{]}; deduce that q=2 (in the contrary case deduce from Ex. 5 that the polynomial \(X^{q^{2}}-a\) would be irreducible). Show further that -a is a square in K (loc. cit.), that -1 is not a square in K, and deduce that \(\Omega=\mathrm{K}(i)\) , where \(i^{2}=-1\) .
\end{enumerate}

\begin{enumerate}
\def\labelenumi{\arabic{enumi})}
\setcounter{enumi}{10}
\item
  Let \(\mathbf{K}\) be a field whose algebraic closure \(\Omega\) is an extension of finite degree \(> 1\) of \(\mathbf{K}\) . Show that \(\Omega = \mathbf{K}(i)\) where \(i^2 = -1\) . (If we had \(\Omega \neq \mathbf{K}(i)\) , show with the help of Galois theory and Sylow's theorem that there would then exist a field \(E\) such that \(\mathbf{K}(i) \subset E \subset \Omega\) and \(\Omega\) is an extension of prime degree of \(E\) ; now apply Ex. 10.)
\item
  Let K be a field of characteristic q, \(\Omega\) an algebraically closed extension of K, x an element of \(\Omega\) such that \(x \notin K\) and M a maximal element of the set of subextensions of \(\Omega\) not containing x (V, p.~156, Ex. 9); suppose further that M is perfect so that \(\mathbf{M}(x)\) is a Galois extension of prime degree p of M (loc. cit.).
\end{enumerate}

\begin{enumerate}
\def\labelenumi{\alph{enumi})}
\item
  Show that either \(p = 2\) and \(\Omega = M(i)\) or for each integer \(r \geqslant 1\) , there exists a single extension of degree \(p^r\) of \(M\) ; this extension \(M_r\) is cyclic over \(M\) , \(\Omega\) is the union of the \(M_r\) and the \(M_r\) are the only extensions of finite degree of \(M\) . (Note that if a \(p\) -group \(\Gamma\) of order \(p^r\) has only one subgroup of order \(p^{r-1}\) , it is necessarily cyclic, arguing by induction on \(r\) and using I, p.~133, Ex. 10, and I, p.~77, Cor. of Prop. 11).
\item
  Suppose that \(\Omega \neq \mathbf{M}(i)\) and that \(p \neq q\) . Show that \(\mathbf{M}\) contains the \(p\) -th roots of unity; we thus have \(\mathbf{M}(x) = \mathbf{M}_1 = \mathbf{M}(\alpha)\) with \(\alpha^p \in \mathbf{M}\) . Show that if \(\alpha\) is the \(p\) -th power of an element of \(\mathbf{M}_1\) , then \(\mathbf{M}_1 = \mathbf{M}(\zeta)\) , where \(\zeta\) is \(p^2\) -th root of unity (if \(\alpha = \beta^p\) for \(\beta \in \mathbf{M}_1\) , consider the minimal polynomial of \(\beta\) in \(\mathbf{M}[\mathbf{X}]\) ).
\item
  Under the hypotheses of \(b\) ) show that if \(\alpha\) is not a \(p\) -th power of an element of \(\mathbf{M}_1\) then \(\mathbf{M}_r\) is the splitting field of the polynomial \(\mathbf{X}^{p^{r-1}} - \alpha\) of \(\mathbf{M}_1[\mathbf{X}]\) (use \(b\) )).
\end{enumerate}

\(d)\) Under the hypotheses of \(b\) ), show that if \(\alpha\) is the \(p\) -th power of an element of \(\mathbf{M}_1\) , there exists \(\gamma \in \mathbf{M}_1\) such that \(\gamma \notin \mathbf{M}\) and that \(\mathbf{M}_2\) is the splitting field of \(\mathbf{X}^p - \gamma\) ; show that \(\gamma\) is not the \(p\) -th power of an element of \(\mathbf{M}_2\) and deduce that \(\mathbf{M}_r\) is the splitting field of the polynomial \(\mathbf{X}^{p^{r-1}} - \gamma\) of \(\mathbf{M}_1[\mathbf{X}]\) .

\begin{enumerate}
\def\labelenumi{\arabic{enumi})}
\setcounter{enumi}{12}
\tightlist
\item
  A field K is called a Moriya field if every separable algebraic extension of finite degree over K is cyclic. * Every finite field is a Moriya field (V, p.~95, Prop. 4). * The field M of Ex. 12 is a Moriya field.
\end{enumerate}

\begin{enumerate}
\def\labelenumi{\alph{enumi})}
\item
  For K to be a Moriya field it is necessary and sufficient that for each integer n \textgreater{} 0 there should exist at most one separable extension of degree n over K. (Use the following fact: if G is a finite group and for each integer \(m \geqslant 1\) dividing the order of G there exists at most one subgroup of G of order m, then G is cyclic; this follows from the same property for p-groups (Ex. 12) and from I, p.~80, Th. 4).
\item
  If K is a Moriya field, show that every algebraic extension E of K is a Moriya field and that if F is an extension of E of degree m over E, then there exists an extension of K of degree m over K.
\end{enumerate}

\begin{enumerate}
\def\labelenumi{\arabic{enumi})}
\setcounter{enumi}{13}
\tightlist
\item
  \begin{enumerate}
  \def\labelenumii{\alph{enumii})}
  \tightlist
  \item
    Let K be a field such that there exist algebraic extensions of K of arbitrarily large finite degree, and that all these extensions have as degree a power of the same prime number p.~Show that under these conditions, for each power \(p^{h}\) ( \(h \geqslant 0\) ) there exists an extension of K of degree \(p^{h}\) (use I, p.~77, Th. 1).
  \end{enumerate}
\end{enumerate}

\begin{enumerate}
\def\labelenumi{\alph{enumi})}
\setcounter{enumi}{1}
\tightlist
\item
  Let p be a prime number and K a field such that there exists an algebraic extension of K of degree divisible by p if \(p \neq 2\) and divisible by 4 if p = 2. Show that there exists an algebraic extension E of K such that the set of degrees of algebraic extensions of E is the set of powers \(p^{h}\) of p.~(If K is perfect or of characteristic \(\neq p\) , let \(K_{1}\) be the perfect closure of K; consider a maximal algebraic extension of \(K_{1}\) , among all those whose elements are of degree not divisible by p over \(K_{1}\) ; use a) and Ex. 11.)
\end{enumerate}

\begin{enumerate}
\def\labelenumi{\arabic{enumi})}
\setcounter{enumi}{14}
\item
  Let K be a Moriya field (Ex. 13) and P the set of prime numbers dividing the degree of an algebraic extension of finite degree of K. Show that the set of degrees of algebraic extensions of finite degree of K is either equal to the set \(N_{P}\) of integers all of whose prime factors are in P or to the subset of \(N_{P}\) consisting of integers not divisible by 4.
\item
  For every prime number \(p\) there exists a Moriya field of characteristic \(p\) such that the set of degrees of algebraic extensions of finite degree of this field is the set \(\mathbf{N}^*\) of integers \(>0^1\) . Show that for every set \(P\) of prime numbers, there exists a Moriya field \(K\) of any characteristic, such that the set of degrees of algebraic extensions of finite degree of \(K\) is the set \(N_P\) (argue as in Ex. 14, b)).
\end{enumerate}

* 17) Let \(\mathbf{P}_n\) be the set of monic polynomials \(\mathbf{X}^n + a_1\mathbf{X}^{n-1} + \cdots + a_n\) of \(\mathbf{Z}[\mathbf{X}]\) whose roots in \(\mathbf{C}\) all have absolute value \(\leqslant 1\) .

\begin{enumerate}
\def\labelenumi{\alph{enumi})}
\item
  Show that the set \(P_{n}\) is finite. Deduce that for every polynomial \(F \in P_{n}\) , if \(F_{h}\) denotes the polynomial whose roots are the h-th powers of the roots of F, then there exist two distinct integers h, k such that \(F_{h} = F_{k}\) .
\item
  Conclude from a) that all the roots of a polynomial \(F \in P_{n}\) are either zero or roots of unity (« Kroneckers's theorem »). *
\end{enumerate}

\begin{enumerate}
\def\labelenumi{\arabic{enumi})}
\setcounter{enumi}{17}
\tightlist
\item
  Let \(p \in \mathbb{N}\) be a prime number, \(C_p\) the cyclotomic extension of level \(p\) of \(\mathbf{Q}\) , \(\zeta \in C_p\) a primitive \(p\) -th root of unity and \(n\) an integer. Write
\end{enumerate}

\[
\mathrm{X} _ {n} (\zeta) = \sum_ {\alpha \in \mathbf {Z} / p \mathbf {Z}} \zeta^ {\alpha^ {n}}.
\]

\begin{enumerate}
\def\labelenumi{\alph{enumi})}
\item
  Calculate \(d = [\mathbf{Q}(\mathbf{X}_n(\zeta)) : \mathbf{Q}]\) . Determine the Galois group of the extension \(\mathbf{Q}(\mathbf{X}_n(\zeta))\) of \(\mathbf{Q}\) and the conjugates of \(\mathbf{X}_n(\zeta)\) .
\item
  Let \(T^{d} + a_{1}T^{d-1} + \cdots + a_{d}\) be the minimal polynomial of \(\mathbf{X}_{n}(\zeta)\) over Q. Show that \(a_{1} = 0\) .
\item
  Suppose that \(n > 1\) and \(p \equiv 1 \pmod{n}\) . Show that \(a_2 = -\frac{n(n - 1)p}{2}\) if \(\frac{p - 1}{n}\) is even and \(a_2 = \frac{np}{2}\) if \(\frac{p - 1}{n}\) is odd.
\item
  Suppose that \(p \equiv 1 \pmod{n}\) . Let \(\beta \in \mathbf{F}_p^*\) be an element whose residue class in \(\mathbf{F}_p^* / \mathbf{F}_p^{*n}\) is a generator. Let \(\mu\) be the number of solutions in \(\mathbf{F}_p^n\) of the equation \(\sum_{i=1}^{n} \beta^{i-1} \mathbf{X}_i^n = 0\) . Show that \(a_n = \frac{(-1)^n p(\mu - p^{n-1})}{p-1}\) .
\item
  Calculate \(a_3\) when \(n = 3\) and \(p = 7, 13, 19\) .
\end{enumerate}

*1 The largest totally ramified extension of the field \(\mathbf{Q}_p\) of \(p\) -adic numbers in its algebraic closure is a field of characteristic 0 with the same property. *

* 19) a) Assume the prime number theorem; if \(\pi(x)\) is the number of primes \(p \leqslant x\) , then \(\pi(x) \sim x / \log x\) as \(x\) tends to \(+\infty\) . Deduce that for every odd integer \(m\) , there exists an increasing sequence of prime numbers \(p_1 < p_2 < \cdots < p_m\) such that \(p_m < p_1 + p_2\) (consider the prime numbers between \(x\) and \(3x/2\) ).

\begin{enumerate}
\def\labelenumi{\alph{enumi})}
\setcounter{enumi}{1}
\tightlist
\item
  With the finite sequence of prime numbers \(p_{j}\) ( \(1 \leq j \leq m\) ) determined as in a), put \(n = p_{1}p_{2} \ldots p_{m}\) . Show that in the cyclotomic polynomial \(\Phi_{n}(X)\) the terms of degree \(< p_{m} + 1\) are the same as those of the polynomial
\end{enumerate}

\[
(1 - \mathbf {X}) ^ {- 1} \prod_ {j = 1} ^ {m} (1 - \mathbf {X} ^ {p _ {j}}).
\]

Deduce that the coefficient of \(\mathbf{X}^{p_m}\) in \(\Phi_n(\mathbf{X})\) is \(1 - m\) (« I. Schur's theorem »). *

* 20) Let \(p\) be a prime number and \(G\) a commutative \(p\) -torsion group (VII, p.~10).\\
a) Let \(G_i \subset G\) be the set of elements \(x \in G\) such that for each \(l\) the equation \(p^{l} y = x\) has a solution. Show that \(G_i\) is a direct factor of \(G\) .

\begin{enumerate}
\def\labelenumi{\alph{enumi})}
\setcounter{enumi}{1}
\item
  Suppose that the equation px = 0 has only a finite number of solutions in G. Show that for each integer l, the kernel \(_{l}G\) of the multiplication by \(p^{l}\) in G is a finite group and that the rank \(r(l)\) of the vector \(F_{p}\) -space \(_{l}G/_{(l-1)}G\) is a decreasing function of l, hence constant for \(l > l_{0}\) , \(l_{0}\) large enough. (Use the theorem of elementary divisors (VII, p.~24). Show that \(G/G_{i}\) is a finite group annihilated by \(p^{l_{0}}\) .)
\item
  Suppose further that \(\mathbf{G} = \mathbf{G}_i\) . Show that the mapping \(l \mapsto r(l)\) is constant. Writing \(\mathbf{H} = (\mathbf{Z}[1/p]/\mathbf{Z})^{(1)}\) , show that for each \(l\) there exists an isomorphism \(_l\mathbf{G} \to {}_l\mathbf{H}\) and that every isomorphism \(_l\mathbf{G} \to {}_l\mathbf{H}\) extends to an isomorphism \((l+1)\mathbf{G} \to {}_{(l+1)}\mathbf{H}\) . Deduce that \(\mathbf{G}\) is isomorphic to \(\mathbf{H}\) .
\item
  Under the hypotheses of b) show that there exists an integer n and an application \(i \to n_i\) of \(N - \{0\}\) into N, with finite support, such that G is isomorphic to \((\mathbf{Z}[1/p]/\mathbf{Z})^n \oplus \bigoplus_{i} (\mathbf{Z}/p^i \mathbf{Z})^{n_i}\) . Show that the integer n and the mapping \(i \mapsto n_i\) are uniquely
\end{enumerate}

determined by the mapping \(l \mapsto \text{Card}(l\mathbf{G})\) . *

* 21) Let G and H be two commutative torsion groups such that for every integer \(n\) , the equation \(nx = 0\) has the same finite number of solutions in G and in H. Show that G is isomorphic to H. (Decompose G and H into direct sums of \(p\) -torsion groups and apply Ex. 20.) Let \(G = \bigoplus_{i \in \mathbf{N}} \mathbf{Z} / p^i\mathbf{Z}\) and \(H = \bigoplus_{i \in \mathbf{N}} \mathbf{Z} / p^{2i}\mathbf{Z}\) . Show that G and H are not isomorphic

and that for each integer \(n\) the cardinal of the set of solutions of the equation \(nx = 0\) is the same in \(G\) and in \(H\) . *

\begin{enumerate}
\def\labelenumi{\arabic{enumi})}
\setcounter{enumi}{21}
\tightlist
\item
  \begin{enumerate}
  \def\labelenumii{\alph{enumii})}
  \tightlist
  \item
    Let k be a field, f a separable monic polynomial in \(k[X]\) , D a splitting extension of \(f, x_{1}, \ldots, x_{n}\) the roots of f in D and G the Galois group of D over k. Assume that G is cyclic. Show that for every element of G to induce an even permutation of the \(x_{i}\) it is necessary and sufficient that the number of irreducible factors of f should be congruent to the degree of f mod 2 (consider a generator \(\sigma\) of G and the decomposition into cycles of the permutation of the \(x_{i}\) induced by \(\sigma\) ).
  \end{enumerate}
\end{enumerate}

\begin{enumerate}
\def\labelenumi{\alph{enumi})}
\setcounter{enumi}{1}
\tightlist
\item
  Suppose further that k is of characteristic \(\neq 2\) . Show that the two above conditions are equivalent to the fact that the discriminant of f should be a square in k (« Stickelberger's theorem »).
\end{enumerate}

\begin{enumerate}
\def\labelenumi{\arabic{enumi})}
\setcounter{enumi}{22}
\tightlist
\item
  Let \(n\) and \(p\) be odd and distinct prime numbers.
\end{enumerate}

\begin{enumerate}
\def\labelenumi{\alph{enumi})}
\item
  Let \(Q \in F_{p}[X]\) be the image of the cyclotomic polynomial \(\Phi_{n}\) . Show that the number i of irreducible factors of Q is even if and only if p is a square mod n.~(Let K be the splitting field of Q. Show that \(\operatorname{Gal}(K/F_{p})\) may be identified with the subgroup of \((\mathbf{F}_{n})^{*}\) generated by the image of p.~Study the parity of the index \(((\mathbf{F}_{n})^{*}: \operatorname{Gal}(K/F_{p}))\) .)
\item
  Show that the discriminant of \(\mathbf{X}^n - 1 \in \mathbf{F}_p[\mathbf{X}]\) is
\end{enumerate}

\[
(- 1) ^ {\frac {n (n - 1)}{2} + n - 1} n ^ {n}.
\]

\begin{enumerate}
\def\labelenumi{\alph{enumi})}
\setcounter{enumi}{2}
\tightlist
\item
  Show that \(i\) is even if and only if \((-1)^{\frac{n(n - 1)}{2}} n\) is a square mod \(p\) (apply Ex. 22 to the polynomial \(\mathbf{X}^n - 1\) ).
\end{enumerate}

\[
\left(\frac {n}{p}\right) = 1
\]

\[
\left(\frac {n}{p}\right) = - 1
\]

\(\left(\frac{n}{p}\right)\left(\frac{p}{n}\right) = (-1)^{\frac{(n - 1)(p - 1)}{4}}\) (« Law of quadratic reciprocity »).
